\documentclass[11pt]{article}
% Shared preamble for the rigor documents (Lamport-style structured proofs).  \input this file after \documentclass.
\usepackage[T1]{fontenc}\usepackage{lmodern}\usepackage{microtype}
\usepackage[margin=1in]{geometry}
\usepackage{amsmath,amssymb,amsthm,mathtools}
\usepackage{xcolor}\usepackage{enumitem}\usepackage{booktabs}\usepackage{graphicx}\usepackage{hyperref}
\usepackage[most]{tcolorbox}
\definecolor{navy}{HTML}{17365D}\definecolor{teal}{HTML}{117A8B}\definecolor{warn}{HTML}{A33A2B}\definecolor{okgreen}{HTML}{2E7D4F}
\hypersetup{colorlinks=true,linkcolor=navy,citecolor=teal,urlcolor=teal}
\theoremstyle{plain}
\newtheorem{theorem}{Theorem}[section]\newtheorem{lemma}[theorem]{Lemma}\newtheorem{proposition}[theorem]{Proposition}\newtheorem{corollary}[theorem]{Corollary}
\theoremstyle{definition}\newtheorem{definition}[theorem]{Definition}\newtheorem{remark}[theorem]{Remark}\newtheorem{claim}[theorem]{Claim}\newtheorem{issue}[theorem]{Issue}
% ---------- Lamport structured proofs ----------
% \lstep{level}{number}{assertion}   -> indented "<level>number. assertion"
% \lproof                             -> "PROOF:" line (start of the sub-proof of the preceding step)
% \lby{justification}                 -> "BY ..." leaf justification for the preceding step
% \lqed{level}{number}                -> QED step of a sub-proof
% keywords: \lassume \lprove \lsuffices \lcase \ldefine \lpick \lnew
\newlength{\lindent}\setlength{\lindent}{1.7em}
\newcommand{\lstep}[3]{\par\smallskip\noindent\hspace*{\dimexpr\numexpr#1-1\relax\lindent\relax}{\color{navy}$\langle#1\rangle#2$.}\ #3}
\newcommand{\lproof}{\par\noindent\hspace*{\lindent}{\scshape Proof:}}
\newcommand{\lby}[1]{\par\noindent\hspace*{\lindent}{\scshape By}\ #1.}
\newcommand{\lqed}[2]{\par\smallskip\noindent\hspace*{\dimexpr\numexpr#1-1\relax\lindent\relax}{\color{navy}$\langle#1\rangle#2$.}\ {\scshape Q.E.D.}}
\newcommand{\lassume}{{\scshape Assume}\ }\newcommand{\lprove}{{\scshape Prove}\ }\newcommand{\lsuffices}{{\scshape Suffices}\ }
\newcommand{\lcase}{{\scshape Case}\ }\newcommand{\ldefine}{{\scshape Define}\ }\newcommand{\lpick}{{\scshape Pick}\ }\newcommand{\lnew}{{\scshape New}\ }
% ---------- verdict boxes ----------
\newtcolorbox{verdictok}{colback=okgreen!8,colframe=okgreen,title={\bfseries Verdict: verified},fonttitle=\small}
\newtcolorbox{verdictgap}{colback=orange!10,colframe=orange!80!black,title={\bfseries Verdict: gap / caveat},fonttitle=\small}
\newtcolorbox{verdictbreak}{colback=warn!10,colframe=warn,title={\bfseries Verdict: proof-breaking issue},fonttitle=\small}
\newtcolorbox{citedbox}[1]{colback=navy!5,colframe=navy!60,title={\bfseries Cited result: #1},fonttitle=\small,breakable}
\setlength{\parindent}{0pt}\setlength{\parskip}{4pt}\allowdisplaybreaks
\newcommand{\cA}{\mathcal A}\newcommand{\cF}{\mathcal F}\newcommand{\cM}{\mathfrak M}\newcommand{\cN}{\mathfrak N}

% extra calligraphic shortcuts (\providecommand: no clash if a document defines its own)
\providecommand{\cB}{\mathcal B}\providecommand{\cC}{\mathcal C}\providecommand{\cD}{\mathcal D}\providecommand{\cE}{\mathcal E}\providecommand{\cG}{\mathcal G}\providecommand{\cH}{\mathcal H}\providecommand{\cI}{\mathcal I}\providecommand{\cJ}{\mathcal J}\providecommand{\cK}{\mathcal K}\providecommand{\cL}{\mathcal L}\providecommand{\cO}{\mathcal O}\providecommand{\cP}{\mathcal P}\providecommand{\cQ}{\mathcal Q}\providecommand{\cR}{\mathcal R}\providecommand{\cS}{\mathcal S}\providecommand{\cT}{\mathcal T}\providecommand{\cU}{\mathcal U}\providecommand{\cV}{\mathcal V}\providecommand{\cW}{\mathcal W}\providecommand{\cX}{\mathcal X}\providecommand{\cY}{\mathcal Y}\providecommand{\cZ}{\mathcal Z}
\newcommand{\Fid}{\operatorname{F}}\newcommand{\Tr}{\operatorname{Tr}}\newcommand{\eps}{\varepsilon}\newcommand{\dd}{\,\mathrm d}


\newcommand{\Wt}{\widetilde W}
\newcommand{\kt}{\widetilde k}
\newcommand{\Stwo}{\mathfrak S_2}
\newcommand{\Pp}{P_+}\newcommand{\Pm}{P_-}
\newcommand{\Om}{\Omega}
\newcommand{\one}{\mathbf 1}
\newcommand{\norm}[1]{\lVert#1\rVert}
\newcommand{\abs}[1]{\lvert#1\rvert}
\newcommand{\ip}[2]{\langle#1,#2\rangle}
\DeclareMathOperator{\dist}{dist}
\DeclareMathOperator{\Var}{Var}

\title{The rate of the remainder in the quadratic law\\
\large Track RR of Phase 2b: $\Phi(\zeta)=f_2\zeta^2\bigl(1+c_3\zeta+O(\zeta^2)\bigr)$,
and the fate of the $\log^{-2}$ claim}
\author{Agent RR}
\date{9 September 2026}

\begin{document}
\maketitle

\begin{abstract}\noindent
Three results.  \textbf{(a)} The Uhlmann--Bogoliubov upper bound of
\texttt{rigor/uhlmann\_upper\_bound.tex} is made quantitative, and the outcome is stronger than
the linear rate the brief asked for: if the extension of the recovery vector field is chosen so that
$s\mapsto\Wt_s$ is a \emph{one-parameter group} (which is possible, and is what the exact integral
representation of Step~4 is made for), then $\norm{V_s}_2^2\le s^2\norm{\Pm D\Pp}_2^2$
\emph{exactly, for all $s\ge0$}, and $\Phi$ obeys the constant-free global bound
$\Phi(\zeta)\le-\frac12\log(1-2f_2\zeta^2)=f_2\zeta^2(1+f_2\zeta^2+O(\zeta^4))$.  The bound is an
\emph{even} function of $\zeta$ and has no cubic term at all; the ``$-0.69\,\zeta$'' read off
\texttt{p1\_bogoliubov.out} is a property of the particular \emph{non}-group family used there,
not of the Uhlmann bound.  Two corollaries:
$c_3\le0$, and the positive $\zeta^2/\log^2(1/\zeta)$ remainder of Note~5 Cor.~4.2 is
\emph{refuted}.  \textbf{(b)} The third-order term of $-\log\Fid$ for quasi-free states is derived
in the eigenbasis of $Q$, checked against the exact determinant formula on random examples, and
combined with the exact defect kernel; the numerical value is $c_3=-0.96(3)$.  \textbf{(c)} The
compressed lower bound of \texttt{rigor/quadratic\_limit.tex} is analysed: the cubic constant
$C^-(\Lambda)$ of the \emph{one-sided} inequality is measured and is bounded, so the optimised
lower rate is \emph{linear}, not $\log^{-2}$.
\medskip\noindent\textbf{After the referee report (agent RRr).}  Three repairs, in
\S\ref{sec:rrr}: the group family is admissible only for $s\ge0$; the hypothesis $f_2^\star=f_2$ is
\emph{removed} (use the group generated by $K=D+ih'$, RRr Prop.~3.1), so
Theorem~\ref{thm:global} and Corollary~\ref{cor:c3sign} are unconditional; and Problem~7.2 of
\texttt{quadratic\_limit.tex} is restated one-sidedly --- its two-sided constant \emph{does} grow
like $e^\Lambda$, as P2 guessed.
\end{abstract}

\tableofcontents

\section{Setting, notation, and the three questions}\label{sec:setting}

All conventions are those of \texttt{rigor/uhlmann\_upper\_bound.tex} (referred to as PA below) and
are recalled only as far as needed.

\begin{itemize}\itemsep2pt
\item $H=L^2(\mathbb R)$, $\Pp=\one_{k>0}$, $\Pm=\one_{k<0}$ (Fourier convention
 $\hat f(k)=\int f(x)e^{-ikx}\dd x$).  For an increasing $C^{1,1}$ diffeomorphism $\phi$ of
 $\mathbb R$, $(W_\phi f)(y)=f(\phi^{-1}(y))\,((\phi^{-1})'(y))^{1/2}$ is unitary.
\item $I=(-a,L)$ with $a=\infty$, $L=1$ in the normal form of PA, so that the cross ratio is
 $\zeta=as/(a+L)=s$.  $I'=(1,\infty)$.  The recovery map is $k_s(x)=x$ on $(-\infty,0]$ and
 $k_s(x)=x/(1+sx)$ on $[0,1]$; it is the time-$s$ flow of the vector field $\chi(x)=x^2$ on
 $[0,1]$, and of $\chi\equiv0$ on $(-\infty,0]$.
\item $\Phi(\zeta):=-\log\Fid_{\cF(I)}(\omega,\widetilde\omega_\zeta)$, the zero-collar recovery
 defect of Note~4/Note~5.  $f_2:=g_B/8=1/(12\pi^2)=0.008443431\ldots$ per chiral component
 (PA Thm.~``A closed''; lower bound \texttt{rigor/quadratic\_limit.tex} Thm.~4.2).
\item $G_\phi:=\Pp-W_\phi^*\Pp W_\phi$ has kernel $\frac i{2\pi}\kappa_\phi$,
 $\kappa_\phi(x,y)=\frac{\sqrt{\phi'(x)\phi'(y)}}{\phi(x)-\phi(y)}-\frac1{x-y}$
 (PA Lemma~2.3), and $E(\phi):=\norm{\Pm W_\phi\Pp}_2^2=\frac12\norm{G_\phi}_2^2
 =\frac1{8\pi^2}\iint\kappa_\phi^2$.
\end{itemize}

The three questions of the brief are: (a) an upper bound $\Phi\le f_2\zeta^2(1+C\zeta)$ with
explicit $C$, and the next coefficient of the Uhlmann bound; (b) the coefficient $c_3$ in
$\Phi=f_2\zeta^2(1+c_3\zeta+O(\zeta^2))$; (c) a lower bound with a rate, from the compressed
curve of \texttt{quadratic\_limit.tex}.  The numerical target is the corrected-basis table of
Note~5 (\texttt{numerics/results\_A2.txt}, \texttt{results\_hp2.txt}):
$\Phi/\zeta^2=0.008376,\;0.008308,\;0.008175$ at $\zeta=\frac1{120},\frac1{60},\frac1{30}$,
i.e.\ deficits $0.80\%$, $1.60\%$, $3.2\%$ relative to $f_2$ --- linear in $\zeta$ with slope
$c_3\approx-0.96$.

\section{(a) The upper bound with a rate}\label{sec:a}

\subsection{Group extensions}

\begin{definition}[Admissible flow extension]\label{def:flow}
A function $\chi:\mathbb R\to\mathbb R$ is an \emph{admissible field} if
(F1) $\chi\in C^{1,1}(\mathbb R)$ with $\Lambda_0:=\norm{\chi'}_\infty<\infty$ and
$\Lambda_1:=\operatorname{Lip}(\chi')<\infty$;
(F2) $\chi(x)=0$ for $x\le0$ and $\chi(x)=x^2$ for $0\le x\le1$;
(F3) $\chi(x)/(1+x^2)\to0$ is \emph{not} required, but $\kappa^{(1)}\in L^2(\mathbb R^2)$ is, where
\begin{equation}\label{eq:k1}
 \kappa^{(1)}(x,y):=\frac1{x-y}\Bigl[\frac{\chi(x)-\chi(y)}{x-y}-\frac{\chi'(x)+\chi'(y)}2\Bigr].
\end{equation}
For an admissible field let $\kt_s$ be its flow, $\partial_s\kt_s(x)=-\chi(\kt_s(x))$,
$\kt_0=\mathrm{id}$, and $\Wt_s:=W_{\kt_s}$, $G_s:=G_{\kt_s}$,
$V_s:=\Pm\Wt_s\Pp$, $E(s):=\norm{V_s}_2^2$.  Put
\begin{equation}\label{eq:f2chi}
 f_2^\chi:=\tfrac12\norm{\Pm D\Pp}_2^2=\frac1{16\pi^2}\iint_{\mathbb R^2}\kappa^{(1)}(x,y)^2\dd x\dd y,
 \qquad D=\chi\partial_x+\tfrac12\chi' .
\end{equation}
\end{definition}

Because of (F2) the flow satisfies $\kt_s|_{(-\infty,0]}=\mathrm{id}$ and
$\kt_s(x)=x/(1+sx)$ on $[0,1]$ \emph{simultaneously for all $s$}: an admissible field produces the
required recovery map on $I$ for the whole one-parameter family at once, \emph{for $s\ge0$}
(for $s<0$ the flow carries $[0,1]$ out of itself and $\kt_s$ is not a recovery map; only
$\zeta=s\ge0$ is ever used).  This is the only property
of the extension used below, and it is exactly the property that the product family
$U_s=\Wt_se^{ish'}$ of PA Definition~4.4 and the Hermite family of
\texttt{numerics/p1\_bogoliubov.py} do \emph{not} have (they are curves, not groups).

\begin{lemma}[The autocorrelation representation]\label{lem:autocorr}
Let $\chi$ be admissible and let
$F_\sigma(x,y):=\sqrt{\kt_\sigma'(x)\kt_\sigma'(y)}\;\kappa^{(1)}(\kt_\sigma(x),\kt_\sigma(y))$.
Then, with $R(u):=\ip{F_0}{F_u}_{L^2(\mathbb R^2)}$,
\begin{equation}\label{eq:autocorr}
 \kappa_{\kt_s}=\int_0^sF_\sigma\dd\sigma\ \text{ in }L^2(\mathbb R^2),\qquad
 \iint\kappa_{\kt_s}^2=\int_0^s\!\!\int_0^s R(\sigma-\tau)\dd\sigma\dd\tau
 =2\int_0^s(s-u)R(u)\dd u,
\end{equation}
and $R$ is real, even, continuous, with $\abs{R(u)}\le R(0)=\iint(\kappa^{(1)})^2$.
\end{lemma}
\lstep{1}{1}{The first identity of \eqref{eq:autocorr}, and $\norm{F_\sigma}_{L^2}=\norm{F_0}_{L^2}$
 for every $\sigma$.}
\lby{PA Lemma~4.2 $\langle1\rangle2$--$\langle1\rangle3$ (the pointwise identity
 $\partial_\sigma\Psi_\sigma=F_\sigma$ with $\Psi_\sigma=\sqrt{\kt'_\sigma\kt'_\sigma}/
 (\kt_\sigma(x)-\kt_\sigma(y))$, $\kappa_{\kt_s}=\Psi_s-\Psi_0$, and the Bochner integral in
 $\Stwo$); the norm identity is the change of variables $X=\kt_\sigma(x)$, $Y=\kt_\sigma(y)$, whose
 Jacobian is exactly $\kt_\sigma'(x)\kt_\sigma'(y)$}
\lstep{1}{2}{$\ip{F_\sigma}{F_\tau}=R(\sigma-\tau)$.}
\lby{$F_\sigma$ is the kernel of $-2\pi i\,\Wt_\sigma^*G^{(1)}\Wt_\sigma$ (PA Lemma~4.2
 $\langle1\rangle3$) and $\Wt_\sigma^*T\Wt_\sigma$ is a unitary conjugation, so
 $\ip{F_\sigma}{F_\tau}=(2\pi)^2\ip{\Wt_\sigma^*G^{(1)}\Wt_\sigma}{\Wt_\tau^*G^{(1)}\Wt_\tau}_{\Stwo}
 =(2\pi)^2\ip{G^{(1)}}{\Wt_{\tau-\sigma}^*G^{(1)}\Wt_{\tau-\sigma}}_{\Stwo}$ by the group law
 $\Wt_\sigma\Wt_\tau=\Wt_{\sigma+\tau}$ (PA Lemma~4.1 $\langle1\rangle1$); the result depends on
 $\tau-\sigma$ only, and is real and symmetric because $\kappa^{(1)}$ is real and symmetric}
\lstep{1}{3}{$R$ is even and $\abs R\le R(0)$; the double integral collapses as stated.}
\lby{$R(-u)=\ip{F_u}{F_0}=R(u)$ by $\langle1\rangle2$ and reality; Cauchy--Schwarz with
 $\langle1\rangle1$; and $\int_0^s\!\int_0^s h(\sigma-\tau)=2\int_0^s(s-u)h(u)\dd u$ for even $h$}
\lqed{1}{4}\lby{$\langle1\rangle1$--$\langle1\rangle3$ and continuity of $u\mapsto R(u)$
 (PA Lemma~4.2 $\langle1\rangle4$: $\sigma\mapsto\Wt^*_\sigma T\Wt_\sigma$ is $\Stwo$-continuous)}

\begin{theorem}[Exact quadratic bound, no error term]\label{thm:exactquad}
For every admissible field $\chi$ and every $s\ge0$ (the inequality and the identities hold for
all real $s$; only $s\ge0$ gives a recovery map, \S\ref{sec:setting}),
\begin{equation}\label{eq:exactquad}
 E(s)=\norm{\Pm\Wt_s\Pp}_2^2\ \le\ 2f_2^\chi\,s^2 ,
\end{equation}
with equality to second order: $E(s)=2f_2^\chi s^2+O(s^4)$.  Moreover $E(-s)=E(s)$, and more
generally $\det(\one-V_s^*V_s)$ is an even function of $s$.
\end{theorem}
\lstep{1}{1}{$E(s)=\frac1{8\pi^2}\iint\kappa_{\kt_s}^2$.}
\lby{PA Prop.~2.4 and the identity $\Tr(G^2)=2\norm{\Pm W\Pp}_2^2$ recorded in
 \texttt{p1\_bogoliubov.out}~(1); equivalently PA's $\langle$PDP$\rangle$ identity with $\kt_s$ in place of $\kt$}
\lstep{1}{2}{$\iint\kappa_{\kt_s}^2\le s^2\iint(\kappa^{(1)})^2=16\pi^2f_2^\chi s^2$.}
\lby{Lemma~\ref{lem:autocorr}: $2\int_0^s(s-u)R(u)\dd u\le R(0)\cdot2\int_0^s(s-u)\dd u=s^2R(0)$}
\lstep{1}{3}{$E(s)=2f_2^\chi s^2+O(s^4)$.}
\lby{$R$ is even, so if $R\in C^2$ near $0$ then $2\int_0^s(s-u)R(u)\dd u
 =s^2R(0)+\frac{s^4}{12}R''(0)+o(s^4)$; without any smoothness of $R$,
 $s^2R(0)\ge\iint\kappa_{\kt_s}^2\ge s^2R(0)-2\int_0^s(s-u)\abs{R(0)-R(u)}\dd u$ and
 $\abs{R(0)-R(u)}=\frac12\norm{F_u-F_0}^2\le\frac12u^2\norm{\kappa^{(2)}}^2$ by
 Lemma~\ref{lem:lip} below, giving the error $\le\frac{s^4}{12}\norm{\kappa^{(2)}}^2$}
\lstep{1}{4}{$E(-s)=E(s)$ and $\det(\one-V^*_{-s}V_{-s})=\det(\one-V_s^*V_s)$.}
\lby{$G_{-s}=\Pp-\Wt_s\Pp\Wt_s^*=-\Wt_sG_s\Wt_s^*$, so $\norm{G_{-s}}_2=\norm{G_s}_2$;
 for the determinant, $V_{-s}=\Pm\Wt_s^*\Pp=(\Pp\Wt_s\Pm)^*$, and in the CS decomposition of the
 unitary $\Wt_s$ with respect to $H=\Pp H\oplus\Pm H$ the two off-diagonal blocks
 $\Pm\Wt_s\Pp$ and $\Pp\Wt_s\Pm$ have the same singular values (the ``sines''), whence the same
 Fredholm determinant}
\lqed{1}{5}\lby{$\langle1\rangle1$--$\langle1\rangle4$}

\subsection{The global upper bound}

\begin{definition}\label{def:f2star}
$f_2^\star:=\inf\{f_2^\chi:\chi\ \text{admissible}\}$.
\end{definition}

\begin{theorem}[Constant-free global upper bound]\label{thm:global}
For every admissible field $\chi$ and every $\zeta\ge0$ with $2f_2^\chi\zeta^2<1$,
\begin{equation}\label{eq:global}
 \Phi(\zeta)\ \le\ -\tfrac12\log\bigl(1-2f_2^\chi\zeta^2\bigr)
 \ =\ f_2^\chi\zeta^2\Bigl(1+f_2^\chi\zeta^2+O(\zeta^4)\Bigr),
\end{equation}
and therefore $\Phi(\zeta)\le-\frac12\log(1-2f_2^\star\zeta^2)$ whenever $2f_2^\star\zeta^2<1$.
In particular the relative remainder of the quadratic law from above is $O(\zeta^2)$, not $O(\zeta)$:
there is no cubic term in the bound.
\end{theorem}
\lstep{1}{1}{$\Phi(\zeta)\le-\frac12\log\det(\one-V_\zeta^*V_\zeta)$ for $\zeta$ with
 $\norm{V_\zeta}<1$.}
\lby{PA Prop.~2.7 (Uhlmann: $\Fid\ge\abs{\ip\Om{\Gamma(\Wt_\zeta)\Om}}$, legitimate because
 $\Wt_\zeta$ implements the recovery map on $\cM=\cF(I)$ by (F2) and acts trivially to the left of
 $0$) and PA Thm.~3.1 (Bogoliubov overlap $=\det(\one-V^*V)^{1/2}$), together with PA Prop.~2.4
 ($\norm{V_\zeta}\le\norm{V_\zeta}_2<\infty$)}
\lstep{1}{2}{$\det(\one-V^*V)\ge1-\norm V_2^2$ whenever $\norm V<1$.}
\lby{$\det(\one-V^*V)=\prod_j(1-\sigma_j^2)$ over the singular values $\sigma_j\in[0,1)$, and
 $\prod_j(1-a_j)\ge1-\sum_ja_j$ for $a_j\in[0,1]$ (induction on the number of factors), with
 $\sum_j\sigma_j^2=\norm V_2^2$}
\lstep{1}{3}{$\norm{V_\zeta}_2^2\le2f_2^\chi\zeta^2$, hence \eqref{eq:global}.}
\lby{Theorem~\ref{thm:exactquad} and $\langle1\rangle1$, $\langle1\rangle2$: $-\log$ is decreasing.
 Note that $\norm{V_\zeta}\le\norm{V_\zeta}_2\le(2f_2^\chi)^{1/2}\zeta<1$ is exactly the standing
 hypothesis, so no smallness assumption beyond $2f_2^\chi\zeta^2<1$ is used}
\lstep{1}{4}{The infimum may be taken inside.}
\lby{$\langle1\rangle3$ holds for each $\chi$; $t\mapsto-\frac12\log(1-2t\zeta^2)$ is continuous and
 increasing on $[0,\frac1{2\zeta^2})$, so the pointwise infimum of the right sides is its value at
 $f_2^\star$}
\lqed{1}{5}\lby{$\langle1\rangle1$--$\langle1\rangle4$}

\begin{proposition}[The value of $f_2^\star$]\label{prop:f2star}
$f_2^\star\ge f_2=1/(12\pi^2)$ unconditionally.  Numerically $f_2^\star=f_2$ to $4\cdot10^{-4}$
relative.
\end{proposition}
\lstep{1}{1}{$f_2^\star\ge f_2$.}
\lby{Theorem~\ref{thm:global} gives $\Phi(\zeta)/\zeta^2\le f_2^\chi+O(\zeta^2)$ for each $\chi$;
 letting $\zeta\downarrow0$ and using $\lim\Phi(\zeta)/\zeta^2=f_2$ (PA Thm.~7.2, ``Theorem A
 closed'': the upper bound is PA Thm.~7.1 and the lower bound is
 \texttt{rigor/quadratic\_limit.tex} Thm.~4.2) yields $f_2\le f_2^\chi$ for every admissible $\chi$}
\lstep{1}{2}{Numerically $f_2^\star\le1.0004\,f_2$.}
\lby{\texttt{numerics/p1\_bogoliubov.out}~(4): minimising the first-order functional
 $Q[g]=\frac1{48\pi^2}\int_0^\infty k^3\abs{\hat g(k)}^2\dd k=\frac1{12}\sum_{n\ge2}(n^3-n)
 \abs{\tilde g_n}^2$ over extensions of the field in a circle-chart Fourier basis of size $M$ gives
 $4.5\,Q_{\min}=0.00846525,\,0.00845812,\,0.00845525,\,0.00845064,\,0.00844939$ for
 $M=30,60,90,140,200$, with Richardson ($1/M$) limit $0.00844649$ against
 $f_2=g_B/8=0.00844343$; each finite $M$ is an admissible smooth field, so
 $f_2^\star\le0.00844939$}
\lqed{1}{3}\lby{$\langle1\rangle1$, $\langle1\rangle2$}

\begin{verdictgap}
$f_2^\star=f_2$ is not proved here.  What is proved is $f_2\le f_2^\star\le0.0084494$
($=1.00071\,f_2$ at $M=200$, $1.00036\,f_2$ after Richardson).  The gap is the same one PA leaves
open in a different guise: PA's Step~5 computes the infimum over the \emph{bounded} commutant
generators $h'\in\mathcal A'$, whereas an extension of the vector field produces the
\emph{unbounded} generator $iD_\eta$, $\eta=\chi_{\mathrm{new}}-\chi$ supported in $I'$.  Both
classes give one-particle--one-hole vectors in the same real subspace $H_{\cM'}$ of PA
Lemma~6.3, and $\dist(a,H_{\cM'})^2=g_B/4$; so $f_2^\star=f_2$ follows as soon as one shows that the
minimiser in $H_{\cM'}$ is a limit of \emph{field} tangents $\Pm(iD_\eta)\Pp$.  The numerics of
\texttt{p1\_bogoliubov.out}~(4) --- a linear least squares in a complete basis, converging like
$1/M$ --- are exactly a minimising sequence of this kind, so I regarded $f_2^\star=f_2$ as
established numerically and open analytically.  \textbf{This verdict is superseded: the gap is
closed in \S\ref{sec:rrr-f2} by enlarging the family from flows to the groups generated by
$K=D+ih'$; $f_2^\star=f_2$ is now a theorem and every conditional statement below is
unconditional.}  Everything below is stated so that the conclusions
survive with $f_2^\star$ in place of $f_2$ wherever it matters.
\end{verdictgap}

\begin{corollary}[Sign of the cubic coefficient; refutation of a positive $\log^{-2}$ remainder]
\label{cor:c3sign}
(Unconditional by \S\ref{sec:rrr-f2}; as originally stated, assume $f_2^\star=f_2$.)  Then
\begin{equation}\label{eq:c3le0}
 \limsup_{\zeta\downarrow0}\frac{\Phi(\zeta)-f_2\zeta^2}{\zeta^3}\ \le\ 0 ,
\end{equation}
so if $\Phi(\zeta)=f_2\zeta^2(1+c_3\zeta+o(\zeta))$ then $c_3\le0$; and there is no remainder term
of the form $+\,b\,\zeta^2/\log^2(1/\zeta)$ with $b>0$: indeed
$\Phi(\zeta)-f_2\zeta^2\le f_2^2\zeta^4/(1-2f_2\zeta^2)$ for all small $\zeta$, whereas
$b\zeta^2/\log^2(1/\zeta)\gg\zeta^4$.  Without the assumption $f_2^\star=f_2$, the same computation
excludes $b=1/120$ (the value of Note~5 Cor.~4.2) for every $\zeta\ge10^{-23}$.
\end{corollary}
\lstep{1}{1}{\eqref{eq:c3le0} and the exclusion of $+b\zeta^2/\log^2(1/\zeta)$.}
\lby{Theorem~\ref{thm:global} with $f_2^\star=f_2$: $\Phi-f_2\zeta^2\le
 -\frac12\log(1-2f_2\zeta^2)-f_2\zeta^2=\sum_{m\ge2}\frac{(2f_2\zeta^2)^m}{2m}
 \le\frac{f_2^2\zeta^4}{1-2f_2\zeta^2}$; divide by $\zeta^3$ and let $\zeta\downarrow0$; and
 $\zeta^{-2}\log^2(1/\zeta)^{-1}$ decays slower than $\zeta^2$}
\lstep{1}{2}{The unconditional statement.}
\lby{With $f_2^\star\le1.00071f_2$ (Prop.~\ref{prop:f2star}), $\Phi-f_2\zeta^2\le
 7.1\cdot10^{-4}f_2\zeta^2+f_2^2\zeta^4/(1-2f_2\zeta^2)$; the Note~5 term
 $\zeta^2/(120\log^2(1/\zeta))$ exceeds $7.1\cdot10^{-4}f_2\zeta^2$ iff
 $\log^2(1/\zeta)<1/(120\cdot7.1\cdot10^{-4}\cdot f_2)=1390$, i.e.\ $\zeta>e^{-37.3}=6\cdot10^{-17}$;
 with the Richardson value $1.00036f_2$ the threshold is $e^{-52.4}=2\cdot10^{-23}$}
\lqed{1}{3}\lby{$\langle1\rangle1$, $\langle1\rangle2$}

\begin{verdictbreak}
Note~5, Corollary~4.2 (A2)--(A3) asserts that the small-$\zeta$ expansion of $\Phi$ carries a
remainder of order $\zeta^2/\log^2(1/\zeta)$, of size $\zeta^2/(120\log^2(1/\zeta))$, coming from
the modular modes $\kappa>\kappa_*\approx2\log(1/\zeta)$ where the second-order (quantum Fisher)
approximation breaks down.  That term is \emph{positive} --- it is $\frac12\int_{\kappa_*}^\infty
(\widetilde q_\kappa-q_\kappa)\dd\kappa>0$ --- and it is $4.3\%$ of $f_2\zeta^2$ at $\zeta=1/120$.
Corollary~\ref{cor:c3sign} rules it out: the Uhlmann--Bogoliubov bound with a group extension
allows a positive remainder of order $\zeta^4$ only.  The numerics agree: the measured remainder is
\emph{negative}, $-0.96\,f_2\zeta^3$.  What is wrong with the heuristic is analysed in
\S\ref{sec:uv}: the modes $\kappa>\kappa_*$ do contribute
$\approx\frac12(\widetilde q_\kappa-q_\kappa)$ each, but this is not an \emph{addition} to
$f_2\zeta^2$: it is (to $4\%$) the \emph{same} quantity that the second-order coefficient already
assigns to those modes, namely $\frac{\zeta^2}8(g_B-g_B^{(\kappa_*)})=\frac{0.513\,\zeta^2}
{8\kappa_*^2}=\frac{\zeta^2}{15.6\,\kappa_*^2}$ against N5's $\frac{\zeta^2}{15\,\kappa_*^2}$
(\texttt{quadratic\_limit.tex} Prop.~7.1 as corrected, and \texttt{numerics/README.md}, which uses
the two interchangeably as the window tail).  The genuine non-analytic remainder is the
\emph{difference} of these two nearly equal quantities, and Corollary~\ref{cor:c3sign} bounds its
positive part by $O(\zeta^4)$.
\end{verdictbreak}


\subsection{The Lipschitz constant, and the quartic coefficient}\label{sec:lip}

The rate the brief asked for --- $\epsilon(s)\le\Lambda_2s$ in PA Prop.~4.5, hence
$\Phi\le f_2\zeta^2(1+C\zeta)$ --- is recorded here because it is what one must use for the
\emph{product} families $U_s=\Wt_se^{ish'}$, which are not groups; for group families
Theorem~\ref{thm:global} supersedes it.  It also gives the quartic coefficient.

\begin{lemma}[Lipschitz continuity of the flow-conjugated first-order kernel]\label{lem:lip}
Let $\chi$ be admissible and assume in addition
\begin{equation}\label{eq:E4}
 \text{(F4)}\qquad \chi\in C^{2,1}\ \text{on}\ [0,\infty)\ \text{and}\ \chi\chi''\ \text{Lipschitz on}
 \ \mathbb R .
\end{equation}
Put $\psi:=\chi\chi'=\frac12(\chi^2)'$, let $\kappa^{(1)}_\psi$ denote the expression \eqref{eq:k1}
with $\psi$ in place of $\chi$, and define
\begin{equation}\label{eq:k2}
 \kappa^{(2)}:=(x-y)\bigl(\kappa^{(1)}\bigr)^2-\kappa^{(1)}_\psi
 +\kappa^{(1)}\Bigl[\frac{\chi(x)-\chi(y)}{x-y}+\frac{\chi'(x)+\chi'(y)}2\Bigr]
 -\frac{(\chi'(x)-\chi'(y))^2}{4(x-y)} .
\end{equation}
Then $\kappa^{(2)}\in L^2(\mathbb R^2)$, $\partial_uF_u(x,y)=\sqrt{\kt_u'(x)\kt_u'(y)}\,
\kappa^{(2)}(\kt_u(x),\kt_u(y))$, and for all $u$
\begin{equation}\label{eq:lip}
 \norm{F_u-F_0}_{L^2}\le\abs u\,\norm{\kappa^{(2)}}_{L^2},\qquad
 \epsilon(u):=\norm{\Wt_u^*G^{(1)}\Wt_u-G^{(1)}}_2\le\frac{\abs u}{2\pi}\norm{\kappa^{(2)}}_{L^2}.
\end{equation}
\end{lemma}
\lstep{1}{1}{With $X=\kt_u(x)$, $Y=\kt_u(y)$, $q:=\frac{\chi(X)-\chi(Y)}{X-Y}$, $a:=\chi'(X)$,
 $b:=\chi'(Y)$ and $A_u:=q-\frac{a+b}2$: $F_u=\Psi_uA_u$ and
 $\partial_uA_u=-(X-Y)\kappa^{(1)}_\psi(X,Y)+q^2-\frac{a^2+b^2}2$.}
\lproof
\lstep{2}{1}{$\partial_uq=\frac{-(\chi'\chi)(X)+(\chi'\chi)(Y)}{X-Y}+q^2$ and
 $\partial_u\bigl(-\frac{a+b}2\bigr)=\frac{(\chi''\chi)(X)+(\chi''\chi)(Y)}2$.}
\lby{$\partial_uX=-\chi(X)$, $\partial_uY=-\chi(Y)$ (definition of the flow) and the quotient rule}
\lstep{2}{2}{$\psi'=(\chi')^2+\chi\chi''$, hence
 $\frac{-\psi(X)+\psi(Y)}{X-Y}+\frac{\psi'(X)+\psi'(Y)}2
 =\partial_uq-q^2+\partial_u(-\frac{a+b}2)+\frac{a^2+b^2}2$.}
\lby{$\langle2\rangle1$ and the definition of $\psi$}
\lqed{2}{3}\lby{$\langle2\rangle2$ rearranged, and
 $\frac{\psi(X)-\psi(Y)}{X-Y}-\frac{\psi'(X)+\psi'(Y)}2=(X-Y)\kappa^{(1)}_\psi(X,Y)$ by \eqref{eq:k1}}
\lstep{1}{2}{$\partial_uF_u=\Psi_u\bigl(A_u^2+\partial_uA_u\bigr)
 =\sqrt{\kt_u'(x)\kt_u'(y)}\,\kappa^{(2)}(X,Y)$ with $\kappa^{(2)}$ as in \eqref{eq:k2}.}
\lby{$F_u=\partial_u\Psi_u=\Psi_uA_u$ (PA Lemma~4.2 $\langle1\rangle2$), so
 $\partial_uF_u=\Psi_u(A_u^2+\partial_uA_u)$; then $\Psi_u=\sqrt{\kt'\kt'}/(X-Y)$,
 $A_u=(X-Y)\kappa^{(1)}(X,Y)$, and the algebraic identity
 $q^2-\frac{a^2+b^2}2=(q-\frac{a+b}2)(q+\frac{a+b}2)-\frac{(a-b)^2}4$}
\lstep{1}{3}{$\kappa^{(2)}\in L^2(\mathbb R^2)$.}
\lproof
\lstep{2}{1}{Near the diagonal all four terms of \eqref{eq:k2} are bounded:
 $\abs{\kappa^{(1)}}\le\Lambda_1$, $\abs{\kappa^{(1)}_\psi}\le\operatorname{Lip}(\psi')$,
 $\abs q,\abs a,\abs b\le\Lambda_0$, $\frac{(a-b)^2}{4\abs{x-y}}\le\frac{\Lambda_1^2}4\abs{x-y}$.}
\lby{PA Lemma~4.2 $\langle1\rangle1$ applied to $\chi$ and to $\psi$; (F4) makes $\psi'$ Lipschitz
 --- this is where the hypothesis is used, and it is \emph{necessary}: if $\chi''$ jumps at a point
 $c$ with $\chi(c)\ne0$ then $\kappa^{(1)}_\psi\sim\text{const}/(x-y)$ on the wedge $y<c<x$ and
 $\kappa^{(2)}\notin L^2$ by a logarithm.  The kink of the recovery field at $x=0$ is harmless
 because $\chi$ and $\chi'$ vanish there to second order, so $\chi\chi''$ is continuous with
 vanishing derivative at $0$; the kink at $x=1$ is removable because the extension is free there}
\lstep{2}{2}{For $y$ outside $\operatorname{supp}\chi$ and $x$ inside, every term is
 $O(\abs{x-y}^{-1})$, hence square integrable.}
\lby{$\chi(y)=\chi'(y)=\psi(y)=\psi'(y)=0$ makes
 $\kappa^{(1)}(x,y)=\frac1{x-y}[\frac{\chi(x)}{x-y}-\frac{\chi'(x)}2]$ and
 $\kappa^{(1)}_\psi$ likewise $O(\abs{x-y}^{-1})$; the four products are then
 $O(\abs{x-y}^{-1})$, and $\int^\infty\abs{x-y}^{-2}\dd y<\infty$}
\lqed{2}{3}\lby{$\langle2\rangle1$, $\langle2\rangle2$ and $\operatorname{supp}\kappa^{(2)}\subseteq
 (K\times\mathbb R)\cup(\mathbb R\times K)$, $K=\operatorname{supp}\chi$}
\lstep{1}{4}{\eqref{eq:lip}.}
\lby{$F_u-F_0=\int_0^u\partial_vF_v\dd v$ pointwise; $\norm{\partial_vF_v}_{L^2}
 =\norm{\kappa^{(2)}}_{L^2}$ for every $v$ by the change of variables of
 Lemma~\ref{lem:autocorr}~$\langle1\rangle1$; then Minkowski's integral inequality.  The second
 form is $\norm{T}_2=\frac1{2\pi}\norm{\text{kernel}}_{L^2}$}
\lqed{1}{5}\lby{$\langle1\rangle1$--$\langle1\rangle4$}

\begin{corollary}[Quartic coefficient; explicit linear rate for product families]\label{cor:quartic}
Under (F1)--(F4),
\begin{equation}\label{eq:quartic}
 2f_2^\chi s^2-\frac{s^4}{96\pi^2}\norm{\kappa^{(2)}}_{L^2}^2\ \le\ E(s)\ \le\ 2f_2^\chi s^2 ,
\end{equation}
and if $R\in C^4$ near $0$ then $E(s)=2f_2^\chi s^2-\frac{s^4}{96\pi^2}\norm{\kappa^{(2)}}^2_{L^2}
+O(s^6)$.  For the product family $U_s=\Wt_se^{ish'}$ of PA Definition~4.4, PA Prop.~4.5
with $\epsilon(s)\le\frac s{2\pi}\norm{\kappa^{(2)}}$ gives
$\bigl|\norm{V^{\mathrm{tot}}_s}_2-s\gamma\bigr|\le(\Lambda_2+C_4)s^2$,
$\gamma=\norm{\Pm(D+ih')\Pp}_2$, $\Lambda_2=\frac1{2\pi}\norm{\kappa^{(2)}}$,
$C_4=C_0\mu+C_0\nu e^\mu+\nu\mu e^\mu$, hence for
$\zeta\le\zeta_0:=\min\{1,\frac1{2\gamma(1+\theta)}\}$ (the bound $\norm{G_s}_2\le C_0s$ on which PA's quadratic expansion, cited here as Prop.~4.5, rests holds for every such $s$ by PA Lemma~5.2, eq.~(21), with its steps $\langle1\rangle1$ and $\langle1\rangle5$ (PA Correction C3, remark), not only in the range $\norm{\mathrm{id}-\kt_s}_\infty\le\frac12$ of PA Prop.~2.5's own proof; Correction~\ref{corr0924:C1}), $\theta:=(\Lambda_2+C_4)/\gamma$,
\begin{equation}\label{eq:linrate}
 \Phi(\zeta)\ \le\ \tfrac12\gamma^2\zeta^2\bigl(1+C\zeta\bigr),\qquad
 C:=2\theta+\theta^2+\tfrac43\gamma^2(1+\theta)^4 .
\end{equation}
\end{corollary}
\lstep{1}{1}{\eqref{eq:quartic}.}
\lby{Lemma~\ref{lem:autocorr} and $R(0)-R(u)=\frac12\norm{F_u-F_0}^2\le\frac{u^2}2
 \norm{\kappa^{(2)}}^2$ (Lemma~\ref{lem:lip}); insert in
 $\iint\kappa^2_{\kt_s}=2\int_0^s(s-u)R(u)\dd u$ and use $2\int_0^s(s-u)\frac{u^2}2\dd u
 =\frac{s^4}{12}$; the upper bound is Theorem~\ref{thm:exactquad}}
\lstep{1}{2}{The Taylor form.}
\lby{$R$ even, so $R(u)=R(0)+\frac{u^2}2R''(0)+O(u^4)$ with
 $R''(0)=-\norm{\kappa^{(2)}}^2$ by $\langle1\rangle1$; then
 $2\int_0^s(s-u)u^2\dd u=\frac{s^4}6$ and $\frac1{8\pi^2}\cdot\frac12\cdot\frac{s^4}6
 =\frac{s^4}{96\pi^2}$}
\lqed{1}{3}\lby{PA Prop.~4.5 with $\epsilon(s)\le\Lambda_2s$; squaring
 $\norm V_2\le\zeta\gamma(1+\theta\zeta)$, then
 $\Phi\le-\frac12\log(1-\norm V_2^2)\le\frac{\norm V_2^2}{2(1-\norm V_2^2)}$ and
 $\frac1{1-v^2}\le1+\frac43v^2$ for $v\le\frac12$}


\subsection{Numerical corroboration, and what \texttt{p1\_bogoliubov.out} really measures}
\label{sec:a-num}

I verified Theorem~\ref{thm:exactquad} directly (script
\texttt{scratchpad/rr\_flow.py}; Cayley chart $x=2\tan\frac\theta2$, staggered $1400\times1400$
periodic trapezoid, flow and its Jacobian by a joint ODE at \texttt{rtol}$=10^{-12}$).  The field is
the admissible one $\chi=0$ ($x\le0$), $\chi=x^2$ ($0\le x\le1$),
$\chi=x^2(1-10t^3+15t^4-6t^5)$, $t=x-1$ ($1\le x\le2$), $\chi=0$ ($x\ge2$), which satisfies
(F1)--(F4).  Its first-order functional, computed independently from the circle Fourier series
$Q=\frac1{12}\sum_{n\ge2}(n^3-n)\abs{\tilde g_n}^2$ (the same routine reproduces
$Q[e^{-x^2/2}]=1/(48\pi)$ to $9$ digits), is $Q=2f_2^\chi=0.146704585$.

\begin{center}\small
\begin{tabular}{r l l l l}\toprule
$s$ & $E(s)$ & $E(-s)$ & $E(s)/s^2$ & $E(s)/(s^2Q)$\\\midrule
$0.05$ & $3.65316338\cdot10^{-4}$ & --- & $0.146126535$ & $0.996060$\\
$0.1$ & $1.44430457\cdot10^{-3}$ & $1.44430455\cdot10^{-3}$ & $0.144430457$ & $0.984499$\\
$0.2$ & $5.52640021\cdot10^{-3}$ & $5.52640008\cdot10^{-3}$ & $0.138160005$ & $0.941757$\\
$0.4$ & $1.90497959\cdot10^{-2}$ & $1.90497949\cdot10^{-2}$ & $0.119061224$ & $0.811571$\\
$0.8$ & $5.26298740\cdot10^{-2}$ & $5.26298680\cdot10^{-2}$ & $0.082234178$ & $0.560543$\\
$1.6$ & $1.11339958\cdot10^{-1}$ & $1.11339896\cdot10^{-1}$ & $0.043492171$ & $0.296461$\\
\bottomrule\end{tabular}\end{center}

Three things are visible.  (i) $E(-s)=E(s)$ to $7$--$8$ digits, the quadrature's own accuracy.
(ii) $E(s)\le Qs^2$ at every $s$, with the ratio decreasing monotonically --- the exact bound
\eqref{eq:exactquad}.  (iii) $\bigl(1-E/(s^2Q)\bigr)/s^2=1.576,\,1.550,\,1.456,\,1.178,\,0.687,
\,0.275$ at $s=0.05,\dots,1.6$: the deficit is \emph{quadratic}, not linear, in $s$
(a linear term would make this column grow like $1/s$), with
$\norm{\kappa^{(2)}}^2/(96\pi^2\cdot2f_2^\chi)=1.58(2)$ from \eqref{eq:quartic}.

\paragraph{The family of \texttt{p1\_bogoliubov.py}.}  That family is \emph{not} a group: the
extension is rebuilt at each $s$ (a cubic Hermite matching $u_s(L),u_s'(L)$, whose data depend
nonlinearly on $s$, plus a correction $\lambda s\sum_ic_i(\dots)$ optimised at first order only).
Its bound therefore may, and does, carry a linear term.  Reading its table~(3) against
$f_2=0.00844343$: $\frac12E/\zeta^2=0.0083923,\,0.0083433,\,0.0082793,\,0.0082463$ at
$\zeta=\frac1{120},\frac1{60},\frac1{30},\frac1{15}$, i.e.\ relative deficits
$0.606\%,\,1.186\%,\,1.944\%,\,2.335\%$, i.e.\ slopes $-0.727,-0.712,-0.583,-0.350$.  So the
brief's ``$-0.69$'' is the small-$\zeta$ slope of \emph{that} curve.  It is consistent with
everything else: $\frac12E\ge\Phi$ requires the slope of the bound to be $\ge c_3=-0.99$, and
$\frac12E/\Phi=1+0.27\zeta$ is exactly $-0.72-(-0.99)=0.27$.

\begin{verdictgap}
There is no contradiction between the even group bound and the odd non-group bound: minimising the
Uhlmann bound over \emph{all} admissible families $U_s$ is a different (and at fixed $\zeta>0$
possibly better) problem than minimising it over one-parameter groups.  What
Theorem~\ref{thm:global} says is that the \emph{group} sub-family already gives, with no
adjustable constant, the sharp quadratic law with an $O(\zeta^2)$ relative error, and that is
enough for Corollary~\ref{cor:c3sign}.  Whether a non-group family can reproduce the true cubic
coefficient $c_3=-0.99$ (its bound would then be sharp to third order) is open; p1's family reaches
$-0.72$.
\end{verdictgap}

\section{(b) The third-order coefficient $c_3$}\label{sec:b}

\subsection{Third order for quasi-free fidelity}

\begin{theorem}[Third-order expansion of $-\log\Fid$]\label{thm:third}
Let $\mathfrak h$ be finite dimensional, $0<Q<1$ diagonal in an orthonormal basis with entries
$q_i$, $c_i:=1-q_i$, $g_i:=q_i/c_i$, $N_{ik}:=q_ic_k+q_kc_i$, and let $D=D^*$ with
$Q_\eps:=Q+\eps D\in(0,1)$ for small $\eps$.  Write $C^2:=1-Q$ and let $T_1,T_2,T_3$ be the
solutions of the Sylvester chain
\begin{equation}\label{eq:sylv}
 GT_1+T_1G=A_1,\qquad GT_2+T_2G=A_2-T_1^2,\qquad GT_3+T_3G=A_3-(T_1T_2+T_2T_1),
\end{equation}
$G=Q(1-Q)^{-1}$, $A_j:=G^{1/2}\,C^{-2}(DC^{-2})^j\,G^{1/2}$; entrywise
$(T_j)_{ik}=(\text{r.h.s.})_{ik}/(g_i+g_k)$.  Put $M_j:=C^2T_j$.  Then
\begin{equation}\label{eq:third}
 -\log\Fid(\omega_Q,\omega_{Q_\eps})=\eps^2\Phi_2+\eps^3\Phi_3+O(\eps^4),\qquad
 \Phi_2=\frac14\sum_{i,k}\frac{\abs{D_{ik}}^2}{N_{ik}},
\end{equation}
\begin{equation}\label{eq:P3}
 \Phi_3=\tfrac16\Tr\bigl[(C^{-2}D)^3\bigr]-\Tr M_3+\Tr(M_1M_2)-\tfrac13\Tr\bigl[M_1^3\bigr],
 \qquad (M_1)_{ik}=\frac{c_i\sqrt{g_ig_k}}{N_{ik}}D_{ik}.
\end{equation}
\end{theorem}
\lstep{1}{1}{$-\log\Fid=-\frac12\log\det(1-Q)-\frac12\log\det(1-Q_\eps)-\log\det(\one+T^{1/2})$
 with $T=G^{1/2}G_\eps G^{1/2}$, $G_\eps=Q_\eps(1-Q_\eps)^{-1}$.}
\lby{\texttt{rigor/check\_note5\_theory.tex} Thm.~2.1 (N5 Thm.~2.1), proved there:
 $\Fid=\det[(1-Q)(1-Q_\eps)]^{1/2}\det[\one+T^{1/2}]$}
\lstep{1}{2}{$-\frac12\log\det(1-Q_\eps)=-\frac12\log\det C^2
 +\frac12\sum_{n\ge1}\frac{\eps^n}n\Tr[(C^{-2}D)^n]$; in particular its $\eps^3$ coefficient is
 $\frac16\Tr[(C^{-2}D)^3]$.}
\lby{$1-Q_\eps=C^2(\one-\eps C^{-2}D)$, $\log\det=\Tr\log$, and
 $-\log(1-t)=\sum_{n\ge1}t^n/n$ (Neumann series, convergent for
 $\eps\norm{C^{-2}D}<1$)}
\lstep{1}{3}{$G_\eps=G+\sum_{n\ge1}\eps^nC^{-2}(DC^{-2})^n$, hence $T=G^2+\sum_{j\ge1}\eps^jA_j$.}
\lby{\texttt{check\_note5\_theory.tex} \S2.3(a) $\langle1\rangle1$ for $n=1,2$ and the same
 computation for general $n$: $(1-Q-\eps D)^{-1}=\sum_n\eps^n C^{-2}(DC^{-2})^n$ and
 $Q_\eps(1-Q_\eps)^{-1}=(1-Q_\eps)^{-1}-\one$, whose $\eps^n$ term is $C^{-2}(DC^{-2})^n$}
\lstep{1}{4}{$T^{1/2}=G+\eps T_1+\eps^2T_2+\eps^3T_3+O(\eps^4)$ with \eqref{eq:sylv}.}
\lby{squaring and matching orders; $G>0$ makes each Sylvester equation uniquely solvable, and in
 the eigenbasis of $Q$ (hence of $G$) it is the entrywise division by $g_i+g_k$}
\lstep{1}{5}{$-\log\det(\one+T^{1/2})=-\log\det(\one+G)-\Tr M+\frac12\Tr M^2-\frac13\Tr M^3
 +O(\eps^4)$, where $M:=(\one+G)^{-1}(T^{1/2}-G)=C^2(T^{1/2}-G)=\eps M_1+\eps^2M_2+\eps^3M_3
 +O(\eps^4)$.}
\lby{$\det(\one+G+\Delta)=\det(\one+G)\det(\one+(\one+G)^{-1}\Delta)$ and $\Tr\log$; and
 $(\one+G)^{-1}=(\one+Q(1-Q)^{-1})^{-1}=1-Q=C^2$}
\lstep{1}{6}{The $\eps^0$ terms cancel and the $\eps^1$ terms cancel.}
\lby{$-\frac12\log\det C^2-\frac12\log\det C^2-\log\det(\one+G)=-\log\det C^2-\log\det C^{-2}=0$;
 at first order, $\frac12\Tr(C^{-2}D)-\Tr M_1$ and
 $\Tr M_1=\sum_i c_i(T_1)_{ii}=\sum_i\frac{c_ig_iD_{ii}}{2q_ic_i}=\frac12\sum_i\frac{D_{ii}}{c_i}
 =\frac12\Tr(C^{-2}D)$, using $(T_1)_{ik}=\sqrt{g_ig_k}D_{ik}/N_{ik}$}
\lstep{1}{7}{The $\eps^2$ coefficient is $\Phi_2=\frac14\sum\abs{D_{ik}}^2/N_{ik}$, and the $\eps^3$
 coefficient is \eqref{eq:P3}.}
\lby{$\langle1\rangle2$ and $\langle1\rangle5$: at order $\eps^2$,
 $\frac14\Tr[(C^{-2}D)^2]-\Tr M_2+\frac12\Tr M_1^2$, which is N5's Prop.~2.3 as re-derived in
 \texttt{check\_note5\_theory.tex} \S2.3(b) (and confirmed numerically below); at order $\eps^3$,
 $\Tr M$ contributes $-\Tr M_3$, $\frac12\Tr M^2$ contributes $\Tr(M_1M_2)$, $-\frac13\Tr M^3$
 contributes $-\frac13\Tr M_1^3$, and $\langle1\rangle2$ contributes $\frac16\Tr[(C^{-2}D)^3]$}
\lqed{1}{8}\lby{$\langle1\rangle1$--$\langle1\rangle7$}

\begin{verdictok}
Verified numerically (\texttt{scratchpad/rr\_third.py}, random $Q$ with $q_i=1/(1+e^{2g_i})$,
$g_i$ standard normal, random Hermitian $D$ with $\norm D_{\mathrm F}=1$, exact $\Fid$ from
Theorem~2.1):

\begin{center}\small
\begin{tabular}{r r r r r r}\toprule
$n$ & $\Phi_2$ (formula) & $\frac14\sum\frac{\abs{D_{ik}}^2}{N_{ik}}$ & $\Phi_3$ &
$\eps$ & $(-\log\Fid-\Phi_2\eps^2)/\eps^3$\\\midrule
$3$ & $0.5099498595$ & $0.5099498595$ & $-0.0644584729$ & $10^{-2}$ & $-0.0610074$\\
 & & & & $3\cdot10^{-3}$ & $-0.0634205$\\
 & & & & $10^{-3}$ & $-0.0641124$\\
$5$ & $0.5774963679$ & $0.5774963679$ & $-0.3677649431$ & $10^{-2}$ & $-0.3568982$\\
 & & & & $3\cdot10^{-3}$ & $-0.3644572$\\
 & & & & $10^{-3}$ & $-0.3666505$\\
\bottomrule\end{tabular}\end{center}
The last column converges to $\Phi_3$ linearly in $\eps$ (the residuals are
$3.45,1.04,0.35\cdot10^{-3}$ and $10.9,3.31,1.11\cdot10^{-3}$, i.e.\ exactly proportional to
$\eps$), and $\Phi_2$ agrees with the known second-order weight to all digits.
\end{verdictok}

\subsection{Evaluation of $c_3$ in the modular frame}\label{sec:c3eval}

\paragraph{Data.}  Write $D(\zeta):=\widetilde Q_\zeta-Q=\zeta\mathcal D_1+\zeta^2\mathcal D_2+O(\zeta^3)$ in the box basis
of \texttt{numerics/compression\_box.py} (modes $e^{ik_j\xi}/\sqrt\Lambda$ on
$[\xi_0-\frac\Lambda2,\xi_0+\frac\Lambda2]$, absolute-basis phase $\xi_0=\log(a/L)$, $a=1$, $L=2$,
so $\zeta=s/3$).  The exact defect matrices $\widehat D(s)=Q-\widetilde Q_s$ are the files
\texttt{Dhat\_exact\_s\{0.025,0.05,0.1,0.2\}\_Lam\{60\_k30,120\_k45\}\_g12.npz}.  I recover
$\mathcal D_1,\mathcal D_2$ by solving the (exactly determined) linear system
$\widehat D(s)=\sum_{k=1}^4s^kA_k$ through the four values of $s$, then $\mathcal D_1=-3A_1$,
$\mathcal D_2=-9A_2$; using only $s\le0.1$ and a cubic fit changes every digit below by $<10^{-5}$.
The compression is the one of \texttt{fidelity\_hp.py}: diagonalise $Q_{\mathrm{box}}$, keep the
eigenvalues $w\in(\varepsilon,1-\varepsilon)$, i.e.\ $\abs\kappa\le\kappa_c=\log\frac{1-\varepsilon}
\varepsilon$, and compress $\mathcal D_1,\mathcal D_2$ with the retained eigenvectors.  Then
\begin{equation}\label{eq:c3split}
 \Phi_P(\zeta)=f_2^{(P)}\zeta^2+\bigl(2\Phi_2^{\mathrm{bil}}[\mathcal D_1,\mathcal D_2]+\Phi_3[\mathcal D_1]\bigr)\zeta^3
 +O(\zeta^4),\qquad
 \Phi_2^{\mathrm{bil}}[A,B]:=\tfrac14\sum_{ik}\frac{\operatorname{Re}(\overline{A_{ik}}B_{ik})}{N_{ik}},
\end{equation}
$f_2^{(P)}=\Phi_2[\mathcal D_1]$, and $c_3^{(P)}:=-C(\kappa_c)/f_2^{(P)}$ with
$C(\kappa_c):=-(2\Phi_2^{\mathrm{bil}}+\Phi_3)$.  Script:
\texttt{scratchpad/rr\_c3c.py} (double precision; the box symbol is accurate to $10^{-14}$, so
$\varepsilon\ge10^{-10}$ is safe).

\begin{center}\small
\begin{tabular}{r r r r r r r}\toprule
$\kappa_c$ & $n$ & $f_2^{(P)}$ & $f_2^{(P)}/f_2$ & $2\Phi_2^{\mathrm{bil}}$ & $\Phi_3$ &
$c_3^{(P)}$\\\midrule
$9.21$ & $31$ & $0.0073889$ & $87.51\%$ & $-0.004114$ & $-0.003999$ & $-1.09805$\\
$11.51$ & $39$ & $0.0078685$ & $93.19\%$ & $-0.003128$ & $-0.005162$ & $-1.05359$\\
$13.82$ & $45$ & $0.0080413$ & $95.24\%$ & $-0.002581$ & $-0.005739$ & $-1.03468$\\
$16.12$ & $53$ & $0.0081715$ & $96.78\%$ & $-0.002068$ & $-0.006266$ & $-1.01993$\\
$18.42$ & $61$ & $0.0082468$ & $97.67\%$ & $-0.001718$ & $-0.006625$ & $-1.01168$\\
$20.72$ & $67$ & $0.0082857$ & $98.13\%$ & $-0.001519$ & $-0.006831$ & $-1.00769$\\
$23.03$ & $73$ & $0.0083128$ & $98.45\%$ & $-0.001364$ & $-0.006990$ & $-1.00502$\\
\bottomrule\end{tabular}\end{center}

($\Lambda=60$, $\kappa_{\max}=30$, $N=91$.  The run with $\Lambda=120$, $\kappa_{\max}=45$,
$N=273$ gives $c_3^{(P)}=-1.09958,-1.05387,-1.03221,-1.01987,-1.01250,-1.00787,-1.00481$ at the
same $\kappa_c$: the \emph{box} dependence is $\le3\cdot10^{-3}$ relative and the whole
$\kappa_c$-dependence is the \emph{window}, exactly as \texttt{numerics/README.md} reports for
$\Phi$ itself.)

\paragraph{Two validations.}  (i) $f_2^{(P)}/f_2=97.67\%$ at $\kappa_c=18.42$ against the
independent prediction $1-\frac{0.513}{8f_2\kappa_c^2}=1-\frac{7.60}{\kappa_c^2}=97.76\%$ of
\texttt{quadratic\_limit.tex} Prop.~7.1 (as corrected); the Richardson ($1/\kappa_c^2$) limit of
the third column is $0.008449$--$0.008465$, i.e.\ $f_2$ to $0.1$--$0.3\%$.
(ii) $C(\kappa_c)=0.008113,0.008290,0.008320,0.008334,0.008343,0.008349,0.008355$ --- a convergent,
\emph{bounded} sequence (this is the input to \S\ref{sec:c} below).  Richardson in $1/\kappa_c^2$
and least squares in $\{1,\kappa_c^{-2},\kappa_c^{-3}\}$ over the last five points give
$C_\infty=0.00837$--$0.00840$.

\begin{equation}\label{eq:c3}
 \boxed{\ c_3=-\frac{C_\infty}{f_2}=-0.99\pm0.01\ }\qquad
 \bigl(C_\infty=0.00838(3),\ f_2=0.00844343\bigr).
\end{equation}
The value $c_3=-1$ is inside the error bar; I have no argument for it.

\paragraph{Consistency with the finite-$\zeta$ table, and the quartic term.}  Comparing with the
\emph{raw window} values of \texttt{results\_A2.txt} (column A, $\varepsilon=10^{-8}$,
$\kappa_c=18.42$, no tail added) at the \emph{same} cutoff, so that no extrapolation enters:
\begin{center}\small
\begin{tabular}{r l l r r}\toprule
$\zeta$ & $\Phi_P(\zeta)$ & $\Phi_P/\zeta^2$ & $(\Phi_P/(f_2^{(P)}\zeta^2)-1)/\zeta$ &
implied $c_4^{(P)}$\\\midrule
$1/120$ & $5.679016\cdot10^{-7}$ & $0.00817778$ & $-1.00440$ & $+0.874$\\
$1/60$ & $2.252719\cdot10^{-6}$ & $0.00810979$ & $-0.99700$ & $+0.881$\\
$1/30$ & $8.863062\cdot10^{-6}$ & $0.00797676$ & $-0.98247$ & $+0.876$\\
$1/15$ & $3.431959\cdot10^{-5}$ & $0.00772191$ & $-0.95478$ & $+0.854$\\
\bottomrule\end{tabular}\end{center}
with $f_2^{(P)}=0.0082468$ and $c_3^{(P)}=-1.01168$ from the table above, and
$c_4^{(P)}:=(c_3^{(P)}-\text{slope})/\zeta$.  The implied quartic coefficient is the same,
$+0.87$, over a factor $8$ in $\zeta$: the perturbative third-order computation and the exact
finite-$\zeta$ fidelities agree completely.  This also explains why fitting the tail-corrected
table against $f_2$ gives the smaller value $-0.96$ quoted in the brief: at
$\zeta=\frac1{120},\frac1{60},\frac1{30}$ the positive quartic term removes
$0.87\zeta=0.7\%,1.4\%,2.9\%$ of the cubic deficit, i.e.\ shifts the apparent slope from $-0.99$
to $-0.98,-0.98,-0.96$ --- and the residual difference is the window-tail model.

\section{(c) The lower bound with a rate}\label{sec:c}

\texttt{rigor/quadratic\_limit.tex} proves $\liminf_{\zeta\to0}\Phi(\zeta)/\zeta^2\ge g_B/8$
(Thm.~4.2) through the compressed curve: for a finite-rank compression $P$ (the box/window),
$\Phi\ge\Phi_P$ and $\Phi_P(\zeta)=\frac{\zeta^2}8g_B^{(P)}+o(\zeta^2)$, with
$0\le g_B-g_B^{(\Lambda)}=0.513\,r\,\Lambda^{-2}$ (Prop.~7.1 as corrected).  Its Problem~7.2 asks for
a cubic remainder \emph{uniform in the cutoff}, guessing $C(\Lambda)\sim e^\Lambda$, ``the natural
one: the third derivative of the finite-dimensional fidelity''.  Only the \emph{one-sided}
inequality is ever used, and only it is at stake here:
\begin{equation}\label{eq:prob72}
 \Phi_P(\zeta)\ \ge\ \frac{\zeta^2}8g_B^{(P)}-C^-(\Lambda)\,\zeta^3 ,\qquad 0\le\zeta\le\zeta_0 .
\end{equation}
(For the two-sided constant P2's guess is \emph{correct}; see \S\ref{sec:rrr-72}.  The two
statements differ because the quartic term of $\Phi_P$ is positive: it helps \eqref{eq:prob72}
while forcing the two-sided constant up.)

\begin{proposition}[Optimising $\Lambda(\zeta)$]\label{prop:optlam}
Granting \eqref{eq:prob72} on $0\le\zeta\le\zeta_0$, with $\Lambda=\kappa_c$ the modular cutoff,
\begin{equation}\label{eq:lowrate}
 \Phi(\zeta)\ \ge\ f_2\zeta^2\Bigl(1-\frac{7.60}{\Lambda^2}-\frac{C^-(\Lambda)}{f_2}\zeta\Bigr)
 \qquad\text{for every }\Lambda,
\end{equation}
and the optimal choice of $\Lambda(\zeta)$ gives the relative lower rate
\begin{center}\small
\begin{tabular}{lll}\toprule
growth of $C^-(\Lambda)$ & optimal $\Lambda(\zeta)$ & relative remainder\\\midrule
$C\,e^{\Lambda}$ & $\log\frac1\zeta-3\log\log\frac1\zeta+O(1)$ &
 $\dfrac{7.60}{\log^2(1/\zeta)}\bigl(1+o(1)\bigr)$\\[4pt]
$c\,\Lambda^m$, $m>0$ & $\bigl(15.2f_2/(mc)\bigr)^{1/(m+2)}\zeta^{-1/(m+2)}$ &
 $O\bigl(\zeta^{2/(m+2)}\bigr)$\\[2pt]
bounded, $C^-(\Lambda)\uparrow C_\infty$ & $\Lambda\to\infty$ at fixed $\zeta$ &
 $\dfrac{C_\infty}{f_2}\,\zeta=0.99\,\zeta$\\
\bottomrule\end{tabular}\end{center}
\end{proposition}
\lstep{1}{1}{\eqref{eq:lowrate}.}
\lby{$\Phi\ge\Phi_P\ge\frac{\zeta^2}8g_B^{(P)}-C^-(\Lambda)\zeta^3$ by \eqref{eq:prob72};
 $g_B^{(P)}\ge g_B^{(\Lambda)}=g_B(1-\frac{0.513}{g_B\Lambda^2})$ (\texttt{quadratic\_limit.tex}
 Thm.~4.2 $\langle1\rangle6$ and Prop.~7.1) and $\frac{0.513}{g_B}=\frac{0.513}{8f_2}=7.60$ for
 $r=1$}
\lstep{1}{2}{The three rows.}
\lby{minimising $h(\Lambda)=7.60\Lambda^{-2}+C^-(\Lambda)\zeta/f_2$: for $C=Ce^\Lambda$,
 $h'=0$ reads $e^\Lambda\Lambda^3=15.2f_2/(C\zeta)$, i.e.\ $\Lambda=\log\frac1\zeta-3\log\Lambda
 +O(1)$, and then both terms are $\Theta(\log^{-2}(1/\zeta))$ with the first equal to
 $7.60\log^{-2}(1/\zeta)(1+o(1))$; for $C=c\Lambda^m$, $h'=0$ reads
 $\Lambda^{m+2}=15.2f_2/(mc\zeta)$ and $h(\Lambda_*)=\Theta(\zeta^{2/(m+2)})$; for bounded $C$ the
 first term can be sent to $0$ at fixed $\zeta$}
\lqed{1}{3}\lby{$\langle1\rangle1$, $\langle1\rangle2$}

\begin{proposition}[$C^-(\Lambda)$ is bounded, numerically]\label{prop:Cbdd}
The exact third Taylor coefficient of $\Phi_P$ at $\zeta=0$ is $-C(\kappa_c)$ with $C(\kappa_c)$ the
column computed in \S\ref{sec:c3eval}: $0.008113$, $0.008290$, $0.008320$, $0.008334$, $0.008343$,
$0.008349$, $0.008355$ at $\kappa_c=9.21,\dots,23.03$ --- monotone increasing, convergent, with
limit $C_\infty=0.00838(3)$, and independent of the box ($\Lambda_{\mathrm{box}}=60$ and $120$
agree to $3\cdot10^{-3}$ relative).  Since $c_3^{(P)}<0<c_4^{(P)}$, the function
$-f_2^{(P)}(c_3^{(P)}+c_4^{(P)}\zeta+\cdots)$ decreases on $[0,\zeta_0]$, so the smallest admissible
$C^-(\Lambda)$ in \eqref{eq:prob72} is attained at $\zeta=0$ and is exactly this third Taylor
coefficient: for the \emph{one-sided} constant, neither exponential nor polynomial growth --- it is
bounded.  (The \emph{two-sided} constant is another matter and does grow like $e^\Lambda$,
\S\ref{sec:rrr-72}.)  Hence the third row of Proposition~\ref{prop:optlam} is the true state of affairs and the
best lower rate obtainable this way is \emph{linear},
\begin{equation}\label{eq:bestlower}
 \Phi(\zeta)\ \ge\ f_2\zeta^2\bigl(1-0.99\,\zeta+O(\zeta^2)\bigr),
\end{equation}
which matches the upper bound $\Phi\le f_2\zeta^2(1+f_2\zeta^2)$ of Theorem~\ref{thm:global} to
the order at which they can match, and is \emph{sharp} (it is the value of $c_3$).
\end{proposition}

\begin{verdictgap}
\eqref{eq:bestlower} is not proved: Proposition~\ref{prop:Cbdd} measures the third
\emph{Taylor coefficient} at $\zeta=0$, whereas \eqref{eq:prob72} needs a bound on the third
derivative on a whole interval, uniformly in $\Lambda$.  What a proof would need, in the notation
of Theorem~\ref{thm:third}:
(i) the off-diagonal exponential decay $\abs{(\mathcal D_1)_{\kappa\kappa'}}\lesssim
e^{-(\kappa+\kappa')/2}$, which is \texttt{check\_note5\_theory.tex} \S5.2 and Cor.~4.2~(A1)
(proved there modulo one contour shift);
(ii) a two-variable power bound $\abs{(\mathcal D_2)_{\kappa\kappa'}}\lesssim
(1+\abs\kappa+\abs{\kappa'})^{-3}$ --- only the \emph{diagonal} $\frac2{15}\zeta^2\kappa^{-3}$ is
proved (N5 Prop.~4.1, \texttt{check\_note5\_theory.tex} Prop.~5.1);
(iii) the cancellation of the two slow tails.  This last point is visible in the table of
\S\ref{sec:c3eval}: $2\Phi_2^{\mathrm{bil}}$ and $\Phi_3$ each converge only like $1/\kappa_c$
($2\Phi_2^{\mathrm{bil}}\approx-31/\kappa_c$; on the diagonal
$\frac{(\mathcal D_1)_{\kappa\kappa}(\mathcal D_2)_{\kappa\kappa}}{4N_{\kappa\kappa}}\approx
\frac{\kappa e^{-\kappa}\cdot\frac2{15}\kappa^{-3}}{8e^{-\kappa}}=\frac1{60}\kappa^{-2}$, whose
tail is $\Theta(1/\kappa_c)$), while their \emph{sum} converges like $1/\kappa_c^2$.  A proof of
uniform boundedness must exhibit that cancellation, not bound the two terms separately.
Bounding them separately gives $C(\Lambda)=O(\Lambda)$, hence by row two of
Proposition~\ref{prop:optlam} the still respectable rate $\Phi\ge f_2\zeta^2(1-O(\zeta^{2/3}))$.
\end{verdictgap}

\begin{verdictok}
For the \emph{one-sided} constant that Thm.~4.2 of \texttt{quadratic\_limit.tex} actually needs, the
exponential guess is not forced: the naive count (a third derivative of a finite-dimensional
fidelity carries $q_{\min}^{-1}\sim e^\Lambda$) is defeated by the $e^{-\kappa}$ decay of the
first-order defect matrix elements, and the measured constant is bounded, so the achievable lower
rate is linear in $\zeta$.  For the two-sided constant of Problem~7.2 as literally stated, P2's
$e^\Lambda$ is right (\S\ref{sec:rrr-72}); an earlier version of this paragraph claimed the
opposite and was wrong.  Correction due to agent RRr.
\end{verdictok}

\section{The ultraviolet heuristic, and the corrected Corollary 4.2}\label{sec:uv}

\subsection{Where the $\log^{-2}$ argument goes wrong}

N5's argument (Cor.~4.2~(A2)--(A3), audited in \texttt{check\_note5\_theory.tex} \S6) is: the
diagonal defect has the power tail $\widetilde q_\kappa-q_\kappa=\frac2{15}\zeta^2\kappa^{-3}$
(N5 Prop.~4.1, proved); for $\kappa>\kappa_*$, where
$\frac2{15}\zeta^2\kappa_*^{-3}=q_{\kappa_*}\approx e^{-\kappa_*}$, i.e.\
$\kappa_*=2\log\frac1\zeta+O(\log\log\frac1\zeta)$, the second-order approximation for such a mode
fails and its contribution saturates at $\frac12(\widetilde q_\kappa-q_\kappa)$; summing,
$\frac12\int_{\kappa_*}^\infty\frac2{15}\frac{\zeta^2}{\kappa^3}\dd\kappa
=\frac{\zeta^2}{30\kappa_*^2}\approx\frac{\zeta^2}{120\log^2(1/\zeta)}$.

Two things are wrong with reading this as a \emph{term} of $\Phi$.

\paragraph{(i) It is not an addition: it is a re-count of what $f_2$ already contains.}
Doubling for $\pm\kappa$, N5's quantity is $\frac{\zeta^2}{15\kappa_*^2}$.  The part of the
second-order coefficient carried by the modes above the same cutoff is
$\frac{\zeta^2}8\bigl(g_B-g_B^{(\kappa_*)}\bigr)=\frac{0.513\,\zeta^2}{8\kappa_*^2}
=\frac{\zeta^2}{15.6\,\kappa_*^2}$ (\texttt{quadratic\_limit.tex} Prop.~7.1 corrected).  These
agree to $4\%$; \texttt{numerics/README.md} uses them interchangeably as \emph{the same} window
tail.  So the modes above $\kappa_*$ contribute about what the quadratic law says they contribute;
the non-analytic remainder is the \emph{difference} of two nearly equal quantities, not either one
of them.

\paragraph{(ii) The mechanism: a UV occupation excess is the shadow of coherent mixing.}
Two modes, $q_1=\epsilon\approx0$ and $q_2=\frac12$, and $\widetilde Q=R^{\!\top}QR$ with $R$ a
rotation by $\theta$.  Then $\widetilde q_1-q_1\approx\frac{\theta^2}2\gg q_1$, so the ``classical''
cost of mode~$1$ would be $\frac12(\widetilde q_1-q_1)=\frac{\theta^2}4$; but the true defect is
$-\log\Fid=\frac14\frac{\abs{D_{12}}^2}{N_{12}}+O(\theta^4)
=\frac14\frac{(\theta/2)^2}{1/2}=\frac{\theta^2}8$, carried entirely by the \emph{off-diagonal}
weight, which is finite because $N_{12}\approx q_2=\frac12$, not $q_1$.  The diagonal excess of a
nearly empty mode is generated by mixing with \emph{low} modes, and its cost is already counted
--- with a benign weight --- in $f_2$.  Counting it again as an independent classical occupation
change over-counts by a factor of order one; the $\log^{-2}$ term is the over-count.

\paragraph{(iii) The data.}  A term $+\frac{\zeta^2}{120\log^2(1/\zeta)}$ would add
$+4.3\%,+5.9\%,+8.5\%$ to $\Phi/(f_2\zeta^2)$ at $\zeta=\frac1{120},\frac1{60},\frac1{30}$
(with the sharper $\kappa_*$ from the transcendental equation, $\kappa_*=20.9,19.1,17.4$:
$+1.8\%,+2.2\%,+2.6\%$).  The measured deviations are $-0.80\%,-1.60\%,-3.18\%$.  Subtracting the
alleged term leaves slopes $-6.1,-4.5,-3.5$ (resp.\ $-3.1,-2.3,-1.7$) --- not remotely constant,
whereas the pure cubic reading gives $-0.96,-0.96,-0.95$, and the cubic-plus-quartic reading of
\S\ref{sec:c3eval} reproduces every digit.  And Corollary~\ref{cor:c3sign} excludes it outright.

\subsection{Corrected statement for Note 5, Corollary 4.2}\label{sec:corrected}

Replace the remainder statement of N5 Cor.~4.2 by:

\begin{quote}\itshape
For the zero-collar recovery of an interval by its adjacent neighbour in the chiral free-fermion
net with $c=r$ components,
\[
 \Phi(\zeta)=f_2\,\zeta^2\bigl(1+c_3\zeta+c_4\zeta^2+O(\zeta^3)\bigr),\qquad
 f_2=\frac r{12\pi^2},\quad c_3=-0.99(1),\quad c_4=+0.9(1),
\]
the expansion being in \emph{integer} powers of $\zeta$; $c_3$ is independent of $r$.  The
remainder after the quadratic term is $O(\zeta^3)$ and \emph{negative}.  In particular there is no
$\zeta^2/\log^2(1/\zeta)$ term: the divergence of the formal fourth-order coefficient
$\sum_\kappa\abs{D^{(2)}_{\kappa\kappa}}^2/N_{\kappa\kappa}$, correctly diagnosed in the earlier
version, is cured by the saturation of the ultraviolet modes, and the saturated contribution of the
modes above $\kappa_*\approx2\log\frac1\zeta$ coincides, to a few percent, with the share
$\frac{\zeta^2}8(g_B-g_B^{(\kappa_*)})=0.513\,r\,\zeta^2/(8\kappa_*^2)$ that the quadratic term
already assigns to them; it is therefore not an additional term of the expansion.  Rigorously:
\[
 f_2\zeta^2\bigl(1-0.99\,\zeta+O(\zeta^2)\bigr)\ \le\ \Phi(\zeta)\ \le\
 -\tfrac12\log\bigl(1-2f_2\zeta^2\bigr)=f_2\zeta^2\bigl(1+f_2\zeta^2+O(\zeta^4)\bigr),
\]
the upper bound proved (Theorem~\ref{thm:global} with \S\ref{sec:rrr-f2}), the lower bound
conditional on a cutoff-uniform cubic remainder (Problem~7.2 of
\texttt{quadratic\_limit.tex}, whose constant is measured to be bounded).
\end{quote}

Also: the last sentence of Cor.~4.2's discussion, ``the fourth-order coefficient is infinite'',
should be kept but relabelled --- the \emph{formal} Taylor coefficient of the quadratic (quantum
Fisher) functional diverges, while $\Phi$ itself has a finite fourth-order coefficient
$c_4f_2\approx+0.0076$.  The two statements are compatible: the divergence sits in an expansion
that is not the expansion of $\Phi$.

\subsection{Summary}
\begin{center}\small
\begin{tabular}{lll}\toprule
& statement & status\\\midrule
upper & $\Phi\le-\frac12\log(1-2f_2^\chi\zeta^2)$, any group extension & proved
 (Thm.~\ref{thm:global})\\
upper & $\Phi\le-\frac12\log(1-2f_2\zeta^2)=f_2\zeta^2(1+f_2\zeta^2+\dots)$ & proved
 (\S\ref{sec:rrr-f2})\\
 & $c_3\le0$; no positive $\zeta^2/\log^2$ remainder & Cor.~\ref{cor:c3sign}\\
value & $c_3=-0.99(1)$, $c_4=+0.9(1)$ & numerical, two routes\\
lower & $\Phi\ge f_2\zeta^2(1-0.99\zeta+O(\zeta^2))$ & conditional on
 $C^-(\Lambda)=O(1)$\\
lower & $\Phi\ge f_2\zeta^2(1-O(\zeta^{2/3}))$ & conditional on $C(\Lambda)=O(\Lambda)$\\
lower & $\Phi\ge f_2\zeta^2(1-7.60/\log^2\frac1\zeta)$ & conditional on $C(\Lambda)=O(e^\Lambda)$\\
\bottomrule\end{tabular}\end{center}

\section{Corrections after the referee report (RRr)}\label{sec:rrr}

Agent RRr's report is \texttt{rigor/referee\_rate\_of\_remainder.tex}.  It endorses the results and
asks for three repairs, applied above and recorded here.

\subsection{(1) The group family is admissible only for $s\ge0$}
The flow of $\chi\ge0$ satisfies $\kt_s([0,1])\subseteq[0,1]$ only forward in time; for $s<0$ the
orbit leaves the interval and $\kt_s$ is not a recovery map.  Definition~\ref{def:flow} ff.\ and
Theorem~\ref{thm:exactquad} now say $s\ge0$.  Nothing downstream changes ($\zeta=s\ge0$ throughout).
RRr also notes correctly that $\Wt_s=e^{sD}$ is \emph{PA's own} family at $h'=0$ (PA Lemma~5.1), not
a new construction: what is new here is only that the group law alone collapses PA's error term.

\subsection{(2) $f_2^\star=f_2$: the hypothesis is removed}\label{sec:rrr-f2}

The restriction to \emph{flows} was gratuitous.  Theorem~\ref{thm:exactquad} uses only the group
law, and the following abstract form makes that explicit.

\begin{lemma}[Group bound; any generator]\label{lem:groupbound}
Let $U_s=e^{sK}$ be a one-parameter unitary group on $H$, $K^*=-K$, with $[\Pm,K]\in\Stwo$, and put
$V_s:=\Pm U_s\Pp$, $G_s:=\Pp-U_s^*\Pp U_s$.  Then, for all $s$,
\begin{equation}\label{eq:groupbound}
 \norm{G_s}_2^2=2\norm{V_s}_2^2\quad\text{(exactly)},\qquad
 \norm{G_s}_2\le\abs s\,\norm{[\Pm,K]}_2,\qquad
 \norm{V_s}_2^2\le s^2\,\norm{\Pm K\Pp}_2^2 ,
\end{equation}
and $\det(\one-V_s^*V_s)$ is an even function of $s$.
\end{lemma}
\lstep{1}{1}{$\norm{G_s}_2^2=2\norm{V_s}_2^2$.}
\lby{$\Tr(G_s^2)=\Tr\Pp+\Tr(U^*\Pp U)-2\Tr(\Pp U^*\Pp U)$ and
 $\Tr(\Pp U^*\Pp U)=\Tr(\Pp U^*U)-\Tr(\Pp U^*\Pm U)=\Tr\Pp-\norm{\Pm U\Pp}_2^2$, all the divergent
 traces cancelling in pairs (rigorously: apply the identity to $\Pi_nG_s\Pi_n$ for an increasing
 sequence of finite-rank projections $\Pi_n\uparrow\one$ and pass to the limit, both sides being monotone)}
\lstep{1}{2}{$G_s=-\int_0^sU_\sigma^*[\Pp,K]U_\sigma\dd\sigma$ as a Bochner integral in $\Stwo$,
 hence $\norm{G_s}_2\le\abs s\norm{[\Pm,K]}_2$.}
\lby{$\frac{\dd}{\dd\sigma}U_\sigma^*\Pp U_\sigma=U_\sigma^*[\Pp,K]U_\sigma$ on a core, the
 integrand is $\Stwo$-continuous with constant norm $\norm{[\Pp,K]}_2=\norm{[\Pm,K]}_2$
 ($\Pp+\Pm=\one$), and the triangle inequality for Bochner integrals}
\lstep{1}{3}{$\norm{[\Pm,K]}_2^2=2\norm{\Pm K\Pp}_2^2$, which with
 $\langle1\rangle1$--$\langle1\rangle2$ gives the third inequality.}
\lby{$K^*=-K$ gives $[\Pm,K]=\Pm K\Pp-\Pp K\Pm$ with $\Pp K\Pm=-(\Pm K\Pp)^*$, and the two
 off-diagonal blocks are $\Stwo$-orthogonal (PA Lemma~5.2 $\langle2\rangle2$)}
\lqed{1}{4}\lby{evenness: for \emph{any} unitary $U$, $\Pm U\Pp$ and $\Pp U\Pm$ have the same
 singular values (CS decomposition; or: $V^*V=\one-A^*A$ and $BB^*=\one-AA^*$ with
 $A=\Pp U\Pp$, and $A^*A$, $AA^*$ are isospectral off $0$), and
 $V_{-s}=\Pm U_s^*\Pp=(\Pp U_s\Pm)^*$.  This is a theorem, not a numerical observation; the table of
 \S\ref{sec:a-num} checks the quadrature, not the mathematics}

\begin{proposition}[RRr Prop.~3.1: $f_2^\star=f_2$, unconditionally]\label{prop:rrr}
For $h'\in\mathcal A'$ (PA Def.~5.3) put $K:=D+ih'$ and $U_s:=e^{sK}$.  Then for every $s\ge0$ the
vector $\Gamma(U_s)^*\Om$ represents $\widetilde\omega_s$ on $\cM=\cF(I)$, and
\begin{equation}\label{eq:rrrbound}
 \Phi(\zeta)\le-\tfrac12\log\bigl(1-2f_2^K\zeta^2\bigr),\quad f_2^K:=\tfrac12\norm{\Pm K\Pp}_2^2,
 \qquad \inf_{h'\in\mathcal A'}f_2^K=\frac{g_B}8=f_2 .
\end{equation}
Hence $\Phi(\zeta)\le-\frac12\log(1-2f_2\zeta^2)$ for $\zeta<(2f_2)^{-1/2}$, with no hypothesis.
\end{proposition}
\lstep{1}{1}{$U_s=\Wt_s c_s$ with the Dyson cocycle
 $c_s=\mathcal T\exp\bigl(i\int_0^sh'_\sigma\dd\sigma\bigr)$,
 $h'_\sigma=\Wt_\sigma^*h'\Wt_\sigma$, and $h'_\sigma=E_{J_\sigma}h'_\sigma E_{J_\sigma}$ with
 $J_\sigma=(\kt_{-\sigma}(1),\infty)\subseteq I'$ for $\sigma\ge0$.}
\lby{RRr Prop.~3.1 $\langle1\rangle1$--$\langle1\rangle2$: $c_s:=\Wt_s^{-1}U_s$ solves
 $\partial_sc_s=(ih'_s)c_s$, $c_0=\one$; and
 $\Wt_\sigma^*E_{I'}\Wt_\sigma=E_{\kt_\sigma^{-1}(I')}$ with $\kt_{-\sigma}(1)\ge1$ because
 $\chi\ge0$ --- \emph{this is where $\sigma\ge0$ is used}}
\lstep{1}{2}{$\Pm h'_\sigma\Pp\in\Stwo$, $\Gamma(c_s)\in\cF(I')$ (twisted), and
 $\Gamma(U_s)^*\Om$ represents $\widetilde\omega_s$ on $\cM$.}
\lby{RRr Prop.~3.1 $\langle1\rangle3$--$\langle1\rangle4$: insert $\one=\Pp+\Pm$ twice in
 $\Pm\Wt_\sigma^*h'\Wt_\sigma\Pp$, each term carrying a factor in $\Stwo$ (PA Def.~5.3 for $\Pm h'\Pp$; for $\Pm\Wt_\sigma\Pp$ and $\Pm\Wt_\sigma^*\Pp=(\Pp\Wt_\sigma\Pm)^*$, PA Cor.~2.7 $\langle1\rangle1$, which after PA's Correction C2 gives both blocks of $\Wt_\sigma$ for every $\sigma$, whereas PA Prop.~2.5 bounds them only where $\sigma\norm\chi_\infty\le\frac12$; Correction~\ref{corr0924:C2});
 then PA Prop.~3.1 and Lemma~3.2 with $\widehat W'$ replaced by $c_s$ (that lemma now assumes both off-diagonal blocks Hilbert--Schmidt, PA Correction C4; for $c_s$ both follow from self-adjointness: $h'_\sigma=h_\sigma'^*$ gives $\Pp h'_\sigma\Pm=(\Pm h'_\sigma\Pp)^*\in\Stwo$, and Duhamel's formula for $\partial_sc_s=(ih'_s)c_s$ as in PA Correction C4(b); Correction~\ref{corr0924:C2}), the commutant factor dropping
 out of $\ip{\Gamma(U_s)^*\Om}{m\Gamma(U_s)^*\Om}$}
\lstep{1}{3}{\eqref{eq:rrrbound}.}
\lby{Lemma~\ref{lem:groupbound} applied to $U_s=e^{sK}$ ($[\Pm,K]\in\Stwo$ by $\langle1\rangle2$ and
 PA Lemma~5.2), then PA Prop.~3.4, PA Thm.~4.1 and $\prod(1-a_j)\ge1-\sum a_j$ exactly as in
 Theorem~\ref{thm:global} $\langle1\rangle1$--$\langle1\rangle3$}
\lqed{1}{4}\lby{$\langle1\rangle3$, PA Prop.~7.6
 ($\inf_{h'\in\mathcal A'}\norm{\Pm(D+ih')\Pp}_2^2=g_B/4$; this is PA's Step~5, proved there) and
 continuity of $t\mapsto-\frac12\log(1-2t\zeta^2)$, so the infimum passes inside}

\begin{verdictok}
Theorem~\ref{thm:global}, Corollary~\ref{cor:c3sign} ($c_3\le0$ and the exclusion of a positive
$\zeta^2/\log^2(1/\zeta)$ remainder) and the corrected N5 Cor.~4.2 of \S\ref{sec:corrected} are now
\emph{unconditional}: the $\zeta\ge6\cdot10^{-17}$ caveat of Corollary~\ref{cor:c3sign}
$\langle1\rangle2$ and the \texttt{verdictgap} after Proposition~\ref{prop:f2star} are void.  Two
steps RRr flags as routine but unwritten: (i) upgrading
$\frac{\dd}{\dd\sigma}U_\sigma^*\Pp U_\sigma=-U_\sigma^*[\Pm,K]U_\sigma$ from weak on the core
$H^1_c\cap D(h')$ to Bochner in $\Stwo$ (Lemma~\ref{lem:groupbound} $\langle1\rangle2$); (ii) PA
Prop.~3.1/Lemma~3.2 are stated for a single localised unitary $\widehat W'$, while $c_s$ is a norm
limit of finite products of such, so continuity of $\Gamma$ on the Shale group is used.  Both are
cheaper than the density statement in $H_{\cM'}$ that the flow-only route required.
\end{verdictok}

\subsection{(3) Problem 7.2: the two-sided constant does grow like $e^\Lambda$}\label{sec:rrr-72}

RRr's objection to the original \texttt{verdictbreak} of \S\ref{sec:c} is correct and is sustained.
A constant valid in the two-sided form
$\bigl|\Phi_P(\zeta)-\frac{\zeta^2}8g_B^{(P)}\bigr|\le C(\Lambda)\zeta^3$ on a \emph{fixed} interval
$[0,\zeta_0]$ must dominate $\zeta_0$ times the quartic Taylor coefficient of $\Phi_P$, and that
coefficient contains the quantum-Fisher weight of the \emph{second}-order defect,
\begin{equation}\label{eq:quarticterm}
 \Phi_2[\mathcal D_2]=\tfrac14\sum_{\abs\kappa,\abs{\kappa'}<\kappa_c}
 \frac{\abs{(\mathcal D_2)_{\kappa\kappa'}}^2}{N_{\kappa\kappa'}},\qquad
 (\mathcal D_2)_{\kappa\kappa}=\tfrac2{15}\abs\kappa^{-3},\quad N_{\kappa\kappa}\approx2e^{-\abs\kappa},
\end{equation}
whose diagonal alone is $\approx e^{\kappa_c}/(225\kappa_c^6)$ --- precisely the divergence N5
Cor.~4.2 diagnoses, an honest derivative at $\zeta=0$, untouched by the saturation that cures it at
finite $\zeta$.  Measured with the matrices of \S\ref{sec:c3eval}
(\texttt{scratchpad/rr\_c4.py}; the full matrix, not only the diagonal):
\begin{center}\small
\begin{tabular}{lrrrrr}\toprule
$\kappa_c$ & $9.21$ & $13.82$ & $18.42$ & $20.72$ & $23.03$\\\midrule
$\Phi_2[\mathcal D_2]$ & $2.81\cdot10^{-3}$ & $1.15\cdot10^{-2}$ & $1.12\cdot10^{-1}$ &
 $3.71\cdot10^{-1}$ & $1.313$\\
$e^{\kappa_c}/(225\kappa_c^6)$ & $7.3\cdot10^{-5}$ & $6.4\cdot10^{-4}$ & $1.14\cdot10^{-2}$ &
 $5.6\cdot10^{-2}$ & $2.98\cdot10^{-1}$\\
\bottomrule\end{tabular}\end{center}
--- a factor $3.3$--$3.5$ per $\Delta\kappa_c=2.3$, i.e.\ exponential, tracking RRr's model up to the
$\kappa_c^{-6}$ prefactor (mine is larger because it includes the off-diagonal weight).  So P2's
guess $C(\Lambda)\sim e^\Lambda$ is right for Problem~7.2 \emph{as literally stated}, and
Proposition~\ref{prop:Cbdd} must be --- and now is --- read as a statement about the one-sided
constant $C^-$ of \eqref{eq:prob72}, which is what Proposition~\ref{prop:optlam} uses and for which
the positive quartic term has the helpful sign.

\paragraph{RRr's suggested decisive test, run.}  Measure the \emph{effective} quartic
$c_4^{(P)}:=\bigl((\Phi_P/(f_2^{(P)}\zeta^2)-1)/\zeta-c_3^{(P)}\bigr)/\zeta$ from the exact window
fidelities at two cutoffs (\texttt{scratchpad/rr\_c4b.py}; $\Phi_P$ raw from
\texttt{results\_A2.txt} column~A at $\kappa_c=18.42$ and \texttt{results\_hp2.txt} at
$\kappa_c=25.33$; $f_2^{(P)},c_3^{(P)}$ from the perturbative code at the same cutoff):
\begin{center}\small
\begin{tabular}{lllrrrr}\toprule
$\kappa_c$ & $f_2^{(P)}$ & $c_3^{(P)}$ & \multicolumn{4}{c}{$c_4^{(P)}$ at
 $\zeta=\frac1{120},\frac1{60},\frac1{30},\frac1{15}$}\\\midrule
$18.42$ & $0.0082468$ & $-1.01168$ & $+0.901$ & $+0.894$ & $+0.880$ & $+0.855$\\
$25.33$ & $0.0083408$ & $-1.00249$ & $+1.384$ & $+1.140$ & $+0.972$ & $+0.877$\\
\bottomrule\end{tabular}\end{center}
Exactly the predicted signature: $c_3^{(P)}$ stays put ($-1.012\to-1.002$, i.e.\ it converges),
while the effective quartic at the widest window is no longer $\zeta$-independent and rises as
$\zeta\downarrow0$ --- the large $\zeta^4$ Taylor term becoming visible once
$\kappa_*(\zeta)\lesssim\kappa_c$.  (The rise is much milder than $\Phi_2[\mathcal D_2]$ alone would
give, so the quartic coefficient has substantial internal cancellation; quantifying it needs the
fourth-order analogue of Theorem~\ref{thm:third} and is left open.)  Neither observation touches
$c_3$, Proposition~\ref{prop:optlam} or Proposition~\ref{prop:Cbdd} in their one-sided reading.

\subsection{(4) Editorial points from the report}

\begin{itemize}\itemsep2pt
\item In the replacement text of \S\ref{sec:corrected}: \emph{theorems} are
 $f_2=r/(12\pi^2)$ (PA Thm.~8.2), the upper bound $\Phi\le-\frac12\log(1-2f_2\zeta^2)$
 (Theorem~\ref{thm:global} with Proposition~\ref{prop:rrr}), the absence of a positive
 $\zeta^2/\log^2(1/\zeta)$ remainder, and $c_3\le0$ given that the limit exists;
 \emph{numerical} are the existence of the limit, $c_3=-0.99(1)$, $c_4=+0.9(1)$ and the reading
 ``integer powers''.  ``$c_3$ is independent of $r$'' is a theorem given the others: by PA
 Thm.~8.1 $\langle1\rangle1$ the $r$-component objects are $r$-fold direct sums, so $\Phi$ and
 $f_2$ are both multiplied by $r$ and every ratio $c_k$ is unchanged.
\item RRr's inherited-hypothesis point: PA Lemma~2.4 and Prop.~2.5 assume $\phi=\mathrm{id}$ outside
 a compact set, whereas the near-optimal extensions are asymptotically affine
 (\texttt{p1\_bogoliubov.out}~(2)).  (F3) is the right substitute and Lemma~\ref{lem:autocorr}
 delivers $\kappa_{\kt_s}\in L^2$, but the $L^2$ bookkeeping of PA Prop.~2.5 and of
 Lemma~\ref{lem:lip} $\langle1\rangle3$ $\langle2\rangle2$ (``for $y$ outside
 $\operatorname{supp}\chi$'') is written for compactly supported fields and must be re-run for
 asymptotic translations.  Recorded as an open bookkeeping item; the kernel computations are
 unaffected.  Note that Proposition~\ref{prop:rrr} makes the whole optimisation over extensions
 unnecessary for the main theorem, so this item is no longer load-bearing.
\item The abstract's ``$\det(\one-V_s^*V_s)$ even in $s$ (no cubic term)'' should not be read as
 licensing a statement about $\Phi$ at $s<0$: the bound majorises $\Phi$ only for $s\ge0$.  The
 conclusion ``no cubic term'' is anyway immediate from the closed form
 $-\frac12\log(1-2f_2^\chi s^2)$.
\end{itemize}

\section{Corrections of 2026-09-24}\label{corr0924:sec}
\sloppy
These corrections implement the knowledge-base card \texttt{c11-extension-implementer} (project
\texttt{cft\_cmi}): its refereed Statement, the DOCUMENT ERROR note recorded from referee
REF-DEP-CFT-B9, and the referee notes of REF-DEP-CFT-B6, B9 and B10 on that card.  PA
(\texttt{rigor/uhlmann\_upper\_bound.tex}) was corrected on the same day (its Corrections C1--C4).
Each changed line was replaced by exactly one line, so line numbers cited elsewhere still point to
the same text.
\begin{enumerate}[label=(C\arabic*),ref=C\arabic*,leftmargin=3em]
\item\label{corr0924:C1} \emph{The bound behind the explicit linear rate}
 (Corollary~\ref{cor:quartic}, l.~376).
 \emph{The document said:} \eqref{eq:linrate} holds for $\zeta\le\zeta_0:=\min\{1,
 \frac1{2\gamma(1+\theta)}\}$, by PA's quadratic expansion (PA Prop.~``Quadratic expansion, with
 error bound'', cited here as Prop.~4.5), without saying where the bound $\norm{G_s}_2\le C_0s$, on
 which that expansion rests, comes from for these $s$.
 \emph{It now says:} the same range, and that this bound holds for every such $s$ by PA Lemma~5.2
 (\texttt{lem:integralrep}), eq.~(21), with its steps $\langle1\rangle1$ and $\langle1\rangle5$.  The statement,
 $\zeta_0$, $\theta$ and $C$ are unchanged.
 \emph{Why:} PA Prop.~2.5's own proof needs $\eta=\norm{\mathrm{id}-\kt_s}_\infty\le\frac12$
 (PA l.~287), and $\eta$ exceeds $\frac12$ inside $[0,1]$: for the standard smooth step of PA,
 $\eta\approx0.6128$ at $s=0.5$ and $\eta\approx0.9538$ at $s=1$
 (\texttt{numerics/kb-checks-2026-09-23/refb6\_flow.py}, output \texttt{refb6\_flow.out}).  So PA
 Prop.~2.5 alone does not cover $\zeta\le\zeta_0$; card \texttt{c11-extension-implementer},
 Statement: ``prop:HS gives an explicit bound when $\norm{\mathrm{id}-\tilde k_s}_\infty\le1/2$
 (l.~286--287) \dots\ its `every $0\le s\le1$' (l.~293) and cor:implementer's `every $s$ in
 $\mathbb R$' use prop:HS beyond its hypothesis''.  The bound itself holds for every real $s$, as
 the referee of PA found on 2026-09-24 (PA Correction C3, remark): eq.~(21) of PA writes
 $G_s=\int_0^s\Wt_\sigma^*G^{(1)}\Wt_\sigma\dd\sigma$ as a Bochner integral in $\Stwo$ for all $s$,
 so $\norm{G_s}_2\le\abs s\,\norm{G^{(1)}}_2$ by unitary invariance; $\norm{G^{(1)}}_2<\infty$ by step
 $\langle1\rangle1$ of that lemma, and $\norm{G^{(1)}}_2=\lim_{s\downarrow0}\norm{G_s}_2/s\le C_0$ by
 its step $\langle1\rangle5$ and PA Prop.~2.5 in its range.  Hence $\norm{G_s}_2\le C_0\abs s$ for
 every $s$, and no restriction on $\zeta$ beyond $\zeta\le\zeta_0$ is needed.
\item\label{corr0924:C2} \emph{Both blocks, for every $\sigma$ and for the cocycle}
 (Proposition~\ref{prop:rrr} $\langle1\rangle2$, l.~868--869).
 \emph{The document said:} each term of $\Pm\Wt_\sigma^*h'\Wt_\sigma\Pp$ carries a factor in
 $\Stwo$ ``(PA Prop.~2.5, Def.~5.3)''; ``then PA Prop.~3.1 and Lemma~3.2 with $\widehat W'$ replaced
 by $c_s$''.
 \emph{It now says:} the factors $\Pm\Wt_\sigma\Pp$ and $\Pm\Wt_\sigma^*\Pp=(\Pp\Wt_\sigma\Pm)^*$
 are Hilbert--Schmidt for every $\sigma$ by PA Cor.~2.7 $\langle1\rangle1$ as corrected by PA's
 Correction C2; PA Prop.~2.5 bounds them only where $\sigma\norm\chi_\infty\le\frac12$.  Both
 off-diagonal blocks of $c_s$ are Hilbert--Schmidt, as PA Lemma~3.2 now requires (PA
 Correction C4(a)): $h'_\sigma$ is self-adjoint, so $\Pp h'_\sigma\Pm=(\Pm h'_\sigma\Pp)^*\in\Stwo$,
 and Duhamel's formula for $\partial_sc_s=(ih'_s)c_s$, $c_0=\one$, gives
 $\Pp c_s\Pm=\int_0^si(\Pp h'_\sigma\Pp\cdot\Pp c_\sigma\Pm+\Pp h'_\sigma\Pm\cdot\Pm c_\sigma\Pm)
 \dd\sigma$, hence $\norm{\Pp c_s\Pm}_2\le e^{s\mu}\int_0^s\norm{\Pp h'_\sigma\Pm}_2\dd\sigma$ with
 $\mu=\norm{h'}$, and likewise for $\Pm c_s\Pp$.  This is the argument of PA Correction C4(b) for
 $e^{ish'}$.
 \emph{Why:} PA Prop.~2.5's explicit bound assumes $\norm{\mathrm{id}-\kt_\sigma}_\infty\le\frac12$
 (Correction~\ref{corr0924:C1}); the Hilbert--Schmidt property for every $\sigma$ comes from
 $\Wt_\sigma=(\Wt_{\sigma/n})^n$.  And Shale--Stinespring needs both off-diagonal blocks, so a
 lemma that uses the implementer of $c_s$ needs both.  Card \texttt{c11-extension-implementer},
 Statement: ``Then $P_-\tilde W_sP_+$ is Hilbert--Schmidt for every $s$.  prop:HS gives an explicit
 bound when $\norm{\mathrm{id}-\tilde k_s}_\infty\le1/2$ \dots\ Every other $s$ follows from
 $\tilde W_s=(\tilde W_{s/n})^n$ \dots\ The other block is $P_+\tilde W_sP_-=(P_-\tilde
 W_{-s}P_+)^*$, Hilbert--Schmidt by the same argument'', and ``The Shale--Stinespring criterion
 needs both off-diagonal blocks; an example is $W'=e^{ish'}$ with $h'$ self-adjoint''; referee
 REF-DEP-CFT-B10: ``Duhamel gives both blocks of $e^{ish'}$ HS''.
\item\label{corr0924:C3} \emph{Known gap (found 2026-09-24, not corrected)}
 (Theorem~\ref{thm:exactquad} $\langle1\rangle1$ and $\langle1\rangle4$, l.~156 and l.~168--171;
 Lemma~\ref{lem:groupbound}, l.~822, 826, 828--832, 842--846).  No proof text was changed for this
 item, and it is not recorded on a knowledge-base card; the same gap is recorded in
 \texttt{rigor/exact\_fidelity\_upper\_half.tex}, Correction C3.
 \emph{What the document says:} in the cosine--sine decomposition of the unitary $\Wt_s$ the two
 off-diagonal blocks $\Pm\Wt_s\Pp$ and $\Pp\Wt_s\Pm$ have the same singular values, whence
 $\det(\one-V_{-s}^*V_{-s})=\det(\one-V_s^*V_s)$.
 \emph{The mechanism:} write a unitary $U$ with respect to $H=\Pp H\oplus\Pm H$ as
 $U=\begin{pmatrix}X&B\\ C&Y\end{pmatrix}$, so that $X=\Pp U\Pp|_{\Pp H}$, $B=\Pp U\Pm$ and
 $C=\Pm U\Pp$.  Unitarity gives $C^*C=\one-X^*X$ and $BB^*=\one-XX^*$.  Since $X^*X$ and $XX^*$ have
 the same spectrum, with multiplicities, away from $0$, $B$ and $C$ have the same singular values
 except for the singular value $1$, whose multiplicity is $\dim\ker X^*$ for $B$ and $\dim\ker X$
 for $C$.  When these kernels are finite-dimensional, the singular values agree if and only if
 $\operatorname{ind}X=\dim\ker X-\dim\ker X^*=0$.  The single smoothly truncated winding of
 \texttt{numerics/kb-checks-2026-09-23/refb9\_winding.py} shows the difference: both of its blocks
 are Hilbert--Schmidt, with squared norms $0.0071$ against $1.0071$ at $R/b=5$, and they differ by
 the index $1$.
 \emph{Affected step:} $\langle1\rangle4$ of Theorem~\ref{thm:exactquad} (the evenness of
 $\det(\one-V_s^*V_s)$ in $s$).  The identity $\Tr(G^2)=2\norm{\Pm W\Pp}_2^2$ quoted in
 $\langle1\rangle1$ (l.~156) is the same statement in squared-norm form, since
 $\Tr(G^2)=\norm{\Pm W\Pp}_2^2+\norm{\Pp W\Pm}_2^2$ when both blocks are Hilbert--Schmidt.  The same
 claim is made for every unitary in Lemma~\ref{lem:groupbound}: its identity
 $\norm{G_s}_2^2=2\norm{V_s}_2^2$ (l.~822; proof l.~828--832, which does not use the group law) and
 its step $\langle1\rangle4$ (l.~842--846).  Through Theorem~\ref{thm:exactquad} $\langle1\rangle1$
 the bound \eqref{eq:exactquad} itself rests on the equality, and with it Theorem~\ref{thm:global}
 and Corollary~\ref{cor:quartic}; Proposition~\ref{prop:rrr} $\langle1\rangle3$ rests on it through
 Lemma~\ref{lem:groupbound}.  Theorem~\ref{thm:global} $\langle1\rangle1$,
 Corollary~\ref{cor:quartic} $\langle1\rangle3$ and Proposition~\ref{prop:rrr} $\langle1\rangle3$
 apply PA Thm.~4.1, which also assumes $\norm{\Pp U\Pm}<1$, with only $\norm V<1$ checked; the
 missing bound follows from $\norm V<1$ by the same equality.  The
 statements hold for unitaries of index zero, $\operatorname{ind}(\Pp\Wt_s\Pp|_{\Pp H})=0$.
 \emph{Suggested repair, not carried out here:} restrict to elements of index zero, for example
 push-forwards by diffeomorphisms such as $\Wt_s$.  This needs its own check, e.g.\ that
 $\operatorname{ind}(\Pp\Wt_s\Pp|_{\Pp H})=0$ along the flow $s\mapsto\Wt_s$; that check is not
 carried out here.
 \emph{Status of the uses (ruling of the referee of these corrections, not yet written as proofs).}
 (i) For the push-forwards $\Wt_s$ (Theorem~\ref{thm:exactquad} $\langle1\rangle1$ and
 $\langle1\rangle4$, Theorem~\ref{thm:global} $\langle1\rangle1$) index zero holds for every $s$:
 $W_\phi$ is a real operator (PA eq.~\texttt{eq:pushforward}), i.e.\ it commutes with complex
 conjugation $\mathcal C$, and $\mathcal C\Pp\mathcal C=\Pm$ (PA eq.~\texttt{eq:hardy}; PA
 Correction C1), so $\Pp\Wt_s\Pm=\mathcal C(\Pm\Wt_s\Pp)\mathcal C$ and the two blocks have the same
 singular values.  This one-line argument is not yet written in any document.
 (ii) For $e^{sK}$, $K=D+ih'$ (Lemma~\ref{lem:groupbound}, Proposition~\ref{prop:rrr}) and the
 product family of Corollary~\ref{cor:quartic}, the documents settle index zero only while
 $\norm{G_s}_2<1$, $G_s=\Pp-U_s^*\Pp U_s$.  For $e^{sK}$ this holds for $\zeta<(4f_2^K)^{-1/2}$ by
 Lemma~\ref{lem:groupbound} $\langle1\rangle2$, since $\norm{G_s}_2\le|s|\,\norm{[\Pm,K]}_2$ and
 $\norm{[\Pm,K]}_2^2=2\norm{\Pm K\Pp}_2^2=4f_2^K$.  Then both blocks have norm $<1$ (they are
 blocks of $[\Pp,U_s]=-U_sG_s$), so $X$ is invertible (PA Thm.~4.1 $\langle1\rangle1$) and
 $\operatorname{ind}X=0$.  The rest needs a standard Fredholm-index continuity argument that no
 document contains, and is open: in Lemma~\ref{lem:groupbound}, which is stated for all $s$, every
 $s$ with $|s|\,\norm{[\Pm,K]}_2\ge1$; in Proposition~\ref{prop:rrr} the range
 $(4f_2^K)^{-1/2}\le\zeta<(2f_2^K)^{-1/2}$; and in Corollary~\ref{cor:quartic} any
 $\zeta\le\zeta_0$ with $\norm{G_\zeta}_2\ge1$.
\end{enumerate}
\end{document}
